\PassOptionsToPackage{unicode=true}{hyperref} % options for packages loaded elsewhere
\PassOptionsToPackage{hyphens}{url}
\PassOptionsToPackage{dvipsnames,svgnames*,x11names*}{xcolor}
%
\documentclass[10pt,]{article}
\usepackage{lmodern}
\usepackage{amssymb,amsmath}
\usepackage{ifxetex,ifluatex}
\usepackage{fixltx2e} % provides \textsubscript
\ifnum 0\ifxetex 1\fi\ifluatex 1\fi=0 % if pdftex
  \usepackage[T1]{fontenc}
  \usepackage[utf8]{inputenc}
  \usepackage{textcomp} % provides euro and other symbols
\else % if luatex or xelatex
  \usepackage{unicode-math}
  \defaultfontfeatures{Ligatures=TeX,Scale=MatchLowercase}
\fi
% use upquote if available, for straight quotes in verbatim environments
\IfFileExists{upquote.sty}{\usepackage{upquote}}{}
% use microtype if available
\IfFileExists{microtype.sty}{%
\usepackage[]{microtype}
\UseMicrotypeSet[protrusion]{basicmath} % disable protrusion for tt fonts
}{}
\IfFileExists{parskip.sty}{%
\usepackage{parskip}
}{% else
\setlength{\parindent}{0pt}
\setlength{\parskip}{6pt plus 2pt minus 1pt}
}
\usepackage{xcolor}
\usepackage{hyperref}
\hypersetup{
            pdftitle={ETM-DVPD Framework --- Current-State Technical Audit},
            colorlinks=true,
            linkcolor=Maroon,
            citecolor=Blue,
            urlcolor=Blue,
            breaklinks=true}
\urlstyle{same}  % don't use monospace font for urls
\usepackage[margin=2.2cm]{geometry}
\usepackage{longtable,booktabs}
% Fix footnotes in tables (requires footnote package)
\IfFileExists{footnote.sty}{\usepackage{footnote}\makesavenoteenv{longtable}}{}
\usepackage{graphicx,grffile}
\makeatletter
\def\maxwidth{\ifdim\Gin@nat@width>\linewidth\linewidth\else\Gin@nat@width\fi}
\def\maxheight{\ifdim\Gin@nat@height>\textheight\textheight\else\Gin@nat@height\fi}
\makeatother
% Scale images if necessary, so that they will not overflow the page
% margins by default, and it is still possible to overwrite the defaults
% using explicit options in \includegraphics[width, height, ...]{}
\setkeys{Gin}{width=\maxwidth,height=\maxheight,keepaspectratio}
\setlength{\emergencystretch}{3em}  % prevent overfull lines
\providecommand{\tightlist}{%
  \setlength{\itemsep}{0pt}\setlength{\parskip}{0pt}}
\setcounter{secnumdepth}{0}
% Redefines (sub)paragraphs to behave more like sections
\ifx\paragraph\undefined\else
\let\oldparagraph\paragraph
\renewcommand{\paragraph}[1]{\oldparagraph{#1}\mbox{}}
\fi
\ifx\subparagraph\undefined\else
\let\oldsubparagraph\subparagraph
\renewcommand{\subparagraph}[1]{\oldsubparagraph{#1}\mbox{}}
\fi

% set default figure placement to htbp
\makeatletter
\def\fps@figure{htbp}
\makeatother


\title{ETM-DVPD Framework --- Current-State Technical Audit}
\providecommand{\subtitle}[1]{}
\subtitle{A self-contained description of the framework as it exists after Phase
0M-R3C}
\date{2026-09-21}

\begin{document}
\maketitle

{
\hypersetup{linkcolor=}
\setcounter{tocdepth}{2}
\tableofcontents
}
\hypertarget{preface-what-this-document-is-and-what-it-is-not}{%
\section{Preface --- what this document is, and what it is
not}\label{preface-what-this-document-is-and-what-it-is-not}}

This is a description of a modelling framework \textbf{as it stands
today}. It is written for a technical auditor who understands power
delivery and packaging and who has never seen this repository. It does
not assume, and does not require, any knowledge of the development
history.

Three things follow from that.

\begin{itemize}
\tightlist
\item
  \textbf{Every claim about the present is generated} from a single
  machine-readable contract,
  \texttt{results/CURRENT\_SCIENTIFIC\_STATUS.yaml::current\_state\_contract}.
  A sentence in this document that contradicts that contract is a build
  failure, not a difference of emphasis. That mechanism exists because
  the previous version of this report was internally consistent with its
  sibling report and consistently stale relative to the actual science
  --- see Section 22.
\item
  \textbf{Missing physics is listed, not hidden.} Sections 20 and 21 say
  what the framework cannot do and which claims are therefore unsafe.
\item
  \textbf{Audited is not implemented.} Where a model was built and
  measured but does not drive the production path, the text says so at
  that point, every time.
\end{itemize}

The companion document \texttt{ETM\_DVPD\_FULL\_TECHNICAL\_REPORT.pdf}
is the chronological record of how the framework reached this state,
including superseded results. Nothing here depends on it.

\newpage

\hypertarget{the-problem-the-framework-solves}{%
\section{1. The problem the framework
solves}\label{the-problem-the-framework-solves}}

For a 1 kW compute die drawing \textbf{1000 A at 1 V}, direct lateral
delivery in the evaluated stack is prohibitively lossy: the IR drop and
the \(I^2R\) loss of a 1 V rail crossing even a few millimetres of
package metal are comparable with the load itself. \textbf{Distributed
vertical power delivery (DVPD)} answers that by converting 48 V down to
1 V immediately above, below or inside the package, so that the
high-current rail travels micrometres rather than millimetres.

That choice is usually made on electrical and thermal grounds. This
framework exists to ask a third question:

\begin{quote}
\textbf{Does package thermo-mechanics change which DVPD architecture you
would choose, or is it only a downstream sign-off check?}
\end{quote}

Answering it requires a chain that runs from an architecture choice all
the way to a local solder-fatigue response, with every link converged
and every parameter traced. The framework is that chain. It is a
\textbf{falsification framework}, not a design tool: each stage was
built so that a negative result would stop the work rather than be
absorbed.

\textbf{The headline answer today.} Mechanics responds strongly and
measurably to the architecture, the ordering of the architectures is
robust, and there is still no admissible mechanical inequality that
would turn that response into a constraint. The framework therefore
RANKS architectures and does not yet CONSTRAIN them.

\hypertarget{the-current-state-in-one-table}{%
\subsection{1.1 The current state, in one
table}\label{the-current-state-in-one-table}}

\begin{longtable}[]{@{}lll@{}}
\toprule
\begin{minipage}[b]{0.30\columnwidth}\raggedright
question\strut
\end{minipage} & \begin{minipage}[b]{0.30\columnwidth}\raggedright
current answer\strut
\end{minipage} & \begin{minipage}[b]{0.30\columnwidth}\raggedright
evidence\strut
\end{minipage}\tabularnewline
\midrule
\endhead
\begin{minipage}[t]{0.30\columnwidth}\raggedright
phase\strut
\end{minipage} & \begin{minipage}[t]{0.30\columnwidth}\raggedright
\textbf{0M-R3C complete}\strut
\end{minipage} & \begin{minipage}[t]{0.30\columnwidth}\raggedright
\texttt{results/phase0mr3c/PHASE0MR3C\_REPARTITION\_AND\_PAPER\_FREEZE.md}\strut
\end{minipage}\tabularnewline
\begin{minipage}[t]{0.30\columnwidth}\raggedright
verdict\strut
\end{minipage} & \begin{minipage}[t]{0.30\columnwidth}\raggedright
\textbf{Phase 0M-R3C: PROCESS-CONDITIONAL REPARTITION FOUND; THE FINAL
Z1/Z2 RESPONSE ORDERING SURVIVES (option B). Phase 0M-R3B counted three
microvia levels from three MATERIAL layers; counting CONDUCTOR PLANES
gives one (DEFECT \#64), which withdraws its
zero-feasible-configurations finding. A permutation of the existing
components, T1C\_REORDERED, makes the 20.408 A switch node an
adjacent-tier RDL transition and drops the worst build-up aspect ratio
from 5.05 to 0.83 at identical component set, per-tier area and 1921 um
stack. One ETM replay on a pair frozen before solving gives
D\_driver\_final = 1.1947 against the historical 1.1943, +0.04 \%.
Neither architecture is PROVEN\_REALIZABLE: geometry closes,
manufacturing does not.}\strut
\end{minipage} & \begin{minipage}[t]{0.30\columnwidth}\raggedright
\texttt{results/phase0mr3a/LS\_TO\_L\_INTERCONNECT\_AUDIT.md}\strut
\end{minipage}\tabularnewline
\begin{minipage}[t]{0.30\columnwidth}\raggedright
ET coupling level\strut
\end{minipage} & \begin{minipage}[t]{0.30\columnwidth}\raggedright
\textbf{ETM-L1: E \textless{}-\textgreater{} T iterated to a fixed
point, T -\textgreater{} M one-way. Intact M -\textgreater{} E is 0.16
\% of R against 22.8 \% from rho(T); M -\textgreater{} T is 1.5 K for a
4 \% TIM bond-line change. Neither is solved back.}\strut
\end{minipage} & \begin{minipage}[t]{0.30\columnwidth}\raggedright
\texttt{results/phase0m/etm\_feedback\_magnitudes.csv}\strut
\end{minipage}\tabularnewline
\begin{minipage}[t]{0.30\columnwidth}\raggedright
mechanical capability\strut
\end{minipage} & \begin{minipage}[t]{0.30\columnwidth}\raggedright
\textbf{Global thermoelastic warpage and stress (M0), plus a convergent
local solder viscoplastic RESPONSE metric (M2) added in Phase 0I and
verified in 0I-R. TWO of eight enumerated mechanisms are addressable,
not one. RDL fracture, delamination, via fatigue, dielectric cracking,
pad cratering and EM-assisted damage remain uncomputed.}\strut
\end{minipage} & \begin{minipage}[t]{0.30\columnwidth}\raggedright
\texttt{results/phase0m/PHYSICS\_PRIORITY\_MATRIX.csv}\strut
\end{minipage}\tabularnewline
\begin{minipage}[t]{0.30\columnwidth}\raggedright
solder metric status\strut
\end{minipage} & \begin{minipage}[t]{0.30\columnwidth}\raggedright
\textbf{AVAILABLE AND CONVERGENT, AND IT IS A RESPONSE. The FE quantity
is a LOCAL FATIGUE-RELEVANT INELASTIC RESPONSE METRIC: accumulated
equivalent inelastic strain per stabilised cycle, volume averaged over
the solder core and matched to the Syed 25 um averaging volume. It is
NOT damage, NOT life and NOT a failure criterion: those require a
sourced damage law, a life model and an admissible allowable
respectively, and none of the three is closed.}\strut
\end{minipage} & \begin{minipage}[t]{0.30\columnwidth}\raggedright
\texttt{results/phase0i/local\_metric\_convergence.csv}\strut
\end{minipage}\tabularnewline
\begin{minipage}[t]{0.30\columnwidth}\raggedright
constitutive status\strut
\end{minipage} & \begin{minipage}[t]{0.30\columnwidth}\raggedright
**CONDITIONAL. The creep card is fitted on Sn3.9Ag0.6Cu, NOT SAC305,
over 7.51-15.92 MPa, while 99.8 \% of the package's accumulated
inelastic strain (dPEEQ*V) is generated outside that window. Two
families can be PARAMETERIZED in the regime; only NORTON can be
INTEGRATED through the production solve. This is the dominant open
blocker.**\strut
\end{minipage} & \begin{minipage}[t]{0.30\columnwidth}\raggedright
\texttt{results/phase0lr/ALLOY\_TRANSFERABILITY\_MATRIX.md}\strut
\end{minipage}\tabularnewline
\begin{minipage}[t]{0.30\columnwidth}\raggedright
life law status\strut
\end{minipage} & \begin{minipage}[t]{0.30\columnwidth}\raggedright
\textbf{AVAILABLE FOR RELATIVE LIFE. The primary sources were obtained
in Phase 0J-R and Syed 2004's author-corrected calibration was
reproduced. HTTP 403 is NO LONGER the blocker; it was, in Phase 0J, and
that statement is historical.}\strut
\end{minipage} & \begin{minipage}[t]{0.30\columnwidth}\raggedright
\texttt{results/phase0jr/PHASE0JR\_PRIMARY\_LIFE\_LAW\_GATE.md}\strut
\end{minipage}\tabularnewline
\begin{minipage}[t]{0.30\columnwidth}\raggedright
required life status\strut
\end{minipage} & \begin{minipage}[t]{0.30\columnwidth}\raggedright
\textbf{OPEN. No field or product requirement exists in this repository
or any supplied source. It is NOT the sole G2 blocker -- Phase 0K
withdrew that framing.}\strut
\end{minipage} & \begin{minipage}[t]{0.30\columnwidth}\raggedright
\texttt{results/phase0k/G2\_DECISION\_REGIONS.md}\strut
\end{minipage}\tabularnewline
\begin{minipage}[t]{0.30\columnwidth}\raggedright
G2 status\strut
\end{minipage} & \begin{minipage}[t]{0.30\columnwidth}\raggedright
\textbf{NOT TESTED, and not testable. THREE independent blockers, any
one sufficient: N\_required is OPEN, absolute N\_f is unvalidated (the
strain and energy formulations disagree 4.1x on this joint), and no
admissible operational mechanical allowable exists (nine candidates
tested against six admissibility conditions, none passed). Closing any
one leaves the other two.}\strut
\end{minipage} & \begin{minipage}[t]{0.30\columnwidth}\raggedright
\texttt{results/phase0h/criterion\_admissibility.csv}\strut
\end{minipage}\tabularnewline
\begin{minipage}[t]{0.30\columnwidth}\raggedright
optimizer status\strut
\end{minipage} & \begin{minipage}[t]{0.30\columnwidth}\raggedright
\textbf{FORBIDDEN. SIX preconditions stand at NO: intra-VR via
technology sourced; process margin on the critical interconnect;
dominant mechanical physics closed; an admissible mechanical criterion
exists; a required life exists; absolute fatigue prediction validated.
Phase 0M-R3C UPGRADED `architecture templates realizable' from NO to
PARTIAL: geometry now closes for a frozen pair, but four manufacturing
inputs remain OPEN so it cannot reach YES.}\strut
\end{minipage} & \begin{minipage}[t]{0.30\columnwidth}\raggedright
\texttt{results/phase0mr3a/final\_open\_items.md}\strut
\end{minipage}\tabularnewline
\begin{minipage}[t]{0.30\columnwidth}\raggedright
cooling status\strut
\end{minipage} & \begin{minipage}[t]{0.30\columnwidth}\raggedright
\textbf{IN PRODUCTION, PHYSICAL, AND NOW REGIME-CHECKED. The production
point at 27 g/s is laminar at Re = 1264.8 and inside the Shah \& London
correlation's validated domain, so every Phase-0M-R conclusion resting
on it stands. Phase 0M-R2 found that the 50 g/s point Phase 0M-R quoted
sits at Re = 2342, above the laminar ceiling of 2300 and below
Gnielinski's floor of 3000, where the two neighbouring correlations
differ by 2.76x: h = 166 000 W/m2K, dp = 1.46 MPa and P\_pump = 74 W are
WITHDRAWN as model outputs. An explicit dispatcher now returns
OUT\_OF\_MODEL\_DOMAIN there rather than extrapolating. Separately, at
27 g/s the thermal entrance length occupies 86 \% of the channel, so the
fully-developed Nusselt number understates h and the production
temperature is CONSERVATIVE. Conjugate CFD validation is reported in
results/phase0mr2/CONJUGATE\_CFD\_VALIDATION.md.}\strut
\end{minipage} & \begin{minipage}[t]{0.30\columnwidth}\raggedright
\texttt{results/phase0mr2/cooling\_regime\_table.csv}\strut
\end{minipage}\tabularnewline
\begin{minipage}[t]{0.30\columnwidth}\raggedright
routing status\strut
\end{minipage} & \begin{minipage}[t]{0.30\columnwidth}\raggedright
\textbf{IN PRODUCTION AND GEOMETRICALLY CORRECTED. Vertical spans are
measured face to face (PHASE0MR3B\_PHYSICAL\_SPANS, used by the
Phase-0M-R3C frozen pair), and via levels are counted from conductor
planes rather than material layers. In the frozen partition ls\_to\_L
and hs\_to\_cf are adjacent-tier RDL transitions at AR 0.083 and the
only build-up via is hs\_to\_ls, 100 um at AR 0.833 carrying 1.435 A.
n\_required is FRAMEWORK PROVISIONING at J\_design = 1e8 A/m2, a
conservative design-side allocation, and is never a manufacturing
qualification. A 3-D DC solve of one physical via column gives 1.16x the
analytical series resistance from land spreading.}\strut
\end{minipage} & \begin{minipage}[t]{0.30\columnwidth}\raggedright
\texttt{results/phase0mr/intravr\_routing\_by\_template.csv}\strut
\end{minipage}\tabularnewline
\begin{minipage}[t]{0.30\columnwidth}\raggedright
crowding status\strut
\end{minipage} & \begin{minipage}[t]{0.30\columnwidth}\raggedright
\textbf{IN PRODUCTION, PER CLASS, CROSS-VALIDATED. The global 1.326
constant is retired from every production path. Each interconnect class
carries its own Delta\_R = R\_3D - R\_1D, solved exactly at its
production geometry and cached. Cross-validated against MOOSE
(independent mesher, formulation and solver) on five geometries: 5/5
ACCEPTED at 0.25-0.73 \%. Delta\_R is a DIFFERENCE against a stated 1-D
reference using identical terminals, NOT an excess: at the nominal TSV
geometry it is NEGATIVE, so the retired constant was inflating the TSV
tier by 32.6 \%.}\strut
\end{minipage} & \begin{minipage}[t]{0.30\columnwidth}\raggedright
\texttt{results/phase0mr/current\_crowding\_validation.csv}\strut
\end{minipage}\tabularnewline
\begin{minipage}[t]{0.30\columnwidth}\raggedright
RDL anisotropy status\strut
\end{minipage} & \begin{minipage}[t]{0.30\columnwidth}\raggedright
\textbf{PROMOTED TO THE PRODUCTION THERMAL DEFAULT. The interface layer
-- the one layer every watt crosses -- was isotropic underfill at k =
0.70 W/mK. Homogenising its own sourced microbump geometry (r = 5 um,
pitch 20 um, h = 10 um) gives kxy = 0.891 and kz = 11.739 W/mK,
i.e.~16.8x more conductive through thickness, within 0.55 \% of the
Voigt bound and 0.61 \% of Maxwell-Garnett. Promotion moves peak
temperature by about 7 K and the driver by about 6.7 \%, and the Z1/Z2
ratio by 0.105 \%: it matters for absolute numbers and changes no
ordering.}\strut
\end{minipage} & \begin{minipage}[t]{0.30\columnwidth}\raggedright
\texttt{results/phase0mr2/rdl\_isotropic\_vs\_anisotropic.csv}\strut
\end{minipage}\tabularnewline
\begin{minipage}[t]{0.30\columnwidth}\raggedright
component partition status\strut
\end{minipage} & \begin{minipage}[t]{0.30\columnwidth}\raggedright
\textbf{GEOMETRICALLY\_REALIZABLE / PROCESS-CONDITIONAL, for both
architectures. The frozen pair is T1C\_REORDERED at a 120 um microvia on
a 150 um pitch: worst aspect ratio 0.833, worst framework routing ratio
0.633, one build-up via level on one net (hs\_to\_ls, 1.435 A). 144
enumerated partitions meet the per-tier area budget and 8 keep both ends
the June-2026 source fixes. PROVEN\_REALIZABLE is NOT available: the
microvia current allowable, the capture-pad diameter, a vendor design
rule and two package clearances are OPEN.}\strut
\end{minipage} & \begin{minipage}[t]{0.30\columnwidth}\raggedright
\texttt{results/phase0mr3a/LS\_TO\_L\_INTERCONNECT\_AUDIT.md}\strut
\end{minipage}\tabularnewline
\begin{minipage}[t]{0.30\columnwidth}\raggedright
next work\strut
\end{minipage} & \begin{minipage}[t]{0.30\columnwidth}\raggedright
\textbf{MANUSCRIPT FREEZE. Phase 0N is NOT authorised and no further
technical phase is scientifically necessary: Phase 0M-R3C found no
fundamental model contradiction and the deciding ordering survived the
entire interconnect correction at +0.04 \%. See
results/PAPER\_FREEZE\_DECISION.md, results/PAPER\_CLAIM\_MATRIX.csv and
results/PAPER\_FIGURE\_TABLE\_PLAN.md. FUTURE WORK, none of which blocks
submission because none bears on a claimed result: high-stress SAC305
characterisation; N\_required; an admissible mechanical allowable and
G2; conjugate CFD; the microvia current allowable and a vendor design
rule; the two remaining package clearances; cohesive delamination, crack
growth and moisture. Items already COMPLETE and not to be relisted:
adopting the physical cooling model into the ETM loop (Phase 0M-R),
putting the RDL into the production thermal model (Phase 0M-R2),
sourcing the interconnect CLASS for ls\_to\_L (Phase 0M-R3B), and the
component repartition (Phase 0M-R3C).}\strut
\end{minipage} & \begin{minipage}[t]{0.30\columnwidth}\raggedright
\texttt{results/phase0mr3a/final\_open\_items.md}\strut
\end{minipage}\tabularnewline
\bottomrule
\end{longtable}

\emph{Contract version 0M-R.1, 2026-09-21. Every current-facing
statement in this report and in the standalone audit report is generated
from this table. A statement that contradicts it is a defect, not a
difference of emphasis.}

\newpage

Report generated \textbar{} 2026-09-21 \textbar{}\\
Repository baseline \textbar{} 2026-09-08 \textbar{}\\
Phase completed \textbar{} 0M-R3C, 2026-09-21 \textbar{}

\textbf{Method.} Twenty-three sequential falsification gates (0B, 0C,
0D, 0D-R, 0E, 0F, 0G, 0H, 0I, 0I-R, 0J, 0J-R, 0K, 0K-R, 0L, 0L-R, 0M,
0M-R, 0M-R2, 0M-R3, 0M-R3A, 0M-R3B, 0M-R3C), each designed so that a
negative result would stop the project rather than be absorbed.

\hypertarget{what-phase-0m-r3c-changed}{%
\subsection{1.2 What Phase 0M-R3C
changed}\label{what-phase-0m-r3c-changed}}

The last planned technical gate, and a paper-freeze decision.

\begin{longtable}[]{@{}ll@{}}
\toprule
\begin{minipage}[b]{0.47\columnwidth}\raggedright
question\strut
\end{minipage} & \begin{minipage}[b]{0.47\columnwidth}\raggedright
answer\strut
\end{minipage}\tabularnewline
\midrule
\endhead
\begin{minipage}[t]{0.47\columnwidth}\raggedright
How many via levels does \texttt{ls\_to\_L} really need?\strut
\end{minipage} & \begin{minipage}[t]{0.47\columnwidth}\raggedright
\textbf{One, not three (DEFECT \#64).} Phase 0M-R3B counted three
MATERIAL layers; \texttt{RDL\_UM} is an interlayer REDISTRIBUTION layer,
which is a conductor plane and therefore a capture land. Counting
conductor planes gives one via level of 100 um\strut
\end{minipage}\tabularnewline
\begin{minipage}[t]{0.47\columnwidth}\raggedright
Does that change the Phase-0M-R3B verdict?\strut
\end{minipage} & \begin{minipage}[t]{0.47\columnwidth}\raggedright
\textbf{Yes.} Its finding that zero configurations satisfy both sourced
bounds is WITHDRAWN. It also relocates the real problem: the binding net
was \texttt{hs\_to\_cf}, crossing the whole 606 um low-side GaN tier at
aspect ratio 5.05\strut
\end{minipage}\tabularnewline
\begin{minipage}[t]{0.47\columnwidth}\raggedright
Can the existing components be repartitioned?\strut
\end{minipage} & \begin{minipage}[t]{0.47\columnwidth}\raggedright
\textbf{Yes, by a permutation.} \texttt{T1C\_REORDERED} swaps gan\_ls
and cap\_fly only. The 20.408 A switch node becomes an adjacent-tier RDL
transition and the worst build-up aspect ratio falls from \textbf{5.05
to 0.83}, at identical component set, per-tier area and 1921 um
stack\strut
\end{minipage}\tabularnewline
\begin{minipage}[t]{0.47\columnwidth}\raggedright
Is it manufacturable?\strut
\end{minipage} & \begin{minipage}[t]{0.47\columnwidth}\raggedright
\textbf{Not demonstrated.} Geometry closes at a source-bounded pitch;
the microvia current allowable, the capture pad, a vendor design rule
and two package clearances are OPEN. The verdict is
\textbf{GEOMETRICALLY\_REALIZABLE / PROCESS-CONDITIONAL}, never
PROVEN\_REALIZABLE\strut
\end{minipage}\tabularnewline
\begin{minipage}[t]{0.47\columnwidth}\raggedright
Does the ordering survive?\strut
\end{minipage} & \begin{minipage}[t]{0.47\columnwidth}\raggedright
\textbf{Yes.} One ETM replay on a pair frozen before solving:
D\_driver\_final = \textbf{1.1947} against the historical 1.1943,
\textbf{+0.03 \%}. ARCH\_Z2 is +8.14 K hotter and carries 16.3 \% less
local fatigue-relevant inelastic response\strut
\end{minipage}\tabularnewline
\bottomrule
\end{longtable}

\textbf{Verdict B --- a process-conditional repartition exists and the
final Z1/Z2 response ordering survives.} No fundamental model
contradiction was found, so framework development stops here. The next
action is manuscript freeze: see
\texttt{results/PAPER\_FREEZE\_DECISION.md}.

\newpage

\hypertarget{superseded-phase-0m-r3b-summary}{%
\subsection{1.3 Superseded Phase-0M-R3B
summary}\label{superseded-phase-0m-r3b-summary}}

Four goals: find out what \texttt{ls\_to\_L} physically is, source the
package clearances, source a real TSV tolerance, and clear the remaining
report-state defects.

\begin{longtable}[]{@{}ll@{}}
\toprule
\begin{minipage}[b]{0.47\columnwidth}\raggedright
question\strut
\end{minipage} & \begin{minipage}[b]{0.47\columnwidth}\raggedright
answer\strut
\end{minipage}\tabularnewline
\midrule
\endhead
\begin{minipage}[t]{0.47\columnwidth}\raggedright
What is \texttt{ls\_to\_L}?\strut
\end{minipage} & \begin{minipage}[t]{0.47\columnwidth}\raggedright
\textbf{Not a 673 um structure.} That span was measured tier mid-plane
to tier mid-plane, through 303 um of GaN die body and 250 um of inductor
body. The build-up a via chain crosses is \textbf{120 um} in three
layers: RDL 10 + flying-capacitor layer 100 + RDL 10 (DEFECT \#62)\strut
\end{minipage}\tabularnewline
\begin{minipage}[t]{0.47\columnwidth}\raggedright
Single via or a chain?\strut
\end{minipage} & \begin{minipage}[t]{0.47\columnwidth}\raggedright
\textbf{A chain} --- one microvia per build-up layer, with a capture
land between each pair. Three levels\strut
\end{minipage}\tabularnewline
\begin{minipage}[t]{0.47\columnwidth}\raggedright
Is the technology sourced?\strut
\end{minipage} & \begin{minipage}[t]{0.47\columnwidth}\raggedright
\textbf{The class is}: stacked build-up microvias through
embedded-component substrate power layers. Two primary sources obtained
in full text. Five of twelve geometric parameters are still OPEN\strut
\end{minipage}\tabularnewline
\begin{minipage}[t]{0.47\columnwidth}\raggedright
Is r = 20 um at 70 um pitch legal?\strut
\end{minipage} & \begin{minipage}[t]{0.47\columnwidth}\raggedright
\textbf{No.} Ostmann 2008 gives min. 150 um I/O pad pitch in production,
100 um at prototype level, and states that below 100 um vialess methods
are used instead. Cauwe 2023 calls three stacked microvia levels a
reliability risk; this gap needs three\strut
\end{minipage}\tabularnewline
\begin{minipage}[t]{0.47\columnwidth}\raggedright
Does routing still pass?\strut
\end{minipage} & \begin{minipage}[t]{0.47\columnwidth}\raggedright
\textbf{No.} The vias available were counted over the 1.80 mm2 die
footprint; the binding number is the 1.61 mm2 left free in the
flying-capacitor layer (DEFECT \#63). At that area the design point
gives \textbf{1.0427 \textgreater{} 1}\strut
\end{minipage}\tabularnewline
\begin{minipage}[t]{0.47\columnwidth}\raggedright
Were the clearances sourced?\strut
\end{minipage} & \begin{minipage}[t]{0.47\columnwidth}\raggedright
\textbf{One of three} moves to SOURCE\_BOUNDED and lands inside its
existing bound, so no corner moves. Two remain OPEN --- and they are no
longer what blocks first\strut
\end{minipage}\tabularnewline
\begin{minipage}[t]{0.47\columnwidth}\raggedright
Is the 32 um TSV margin process-qualified?\strut
\end{minipage} & \begin{minipage}[t]{0.47\columnwidth}\raggedright
\textbf{No.} No manufacturing tolerance could be sourced. 32 um absorbs
3.23 \% symmetric tolerance before AR = 10 is violated; it is
\textbf{parametrically margined, not process-qualified}\strut
\end{minipage}\tabularnewline
\bottomrule
\end{longtable}

\textbf{Verdict D --- the physical \texttt{ls\_to\_L} topology,
evaluated against sourced substrate-via rules, makes the frozen Z1/Z2
pair NOT\_REALIZABLE under the current component topology.} Both
architectures share one converter cell, so it applies to both. The
deciding mechanical ordering is untouched and is not what this gate
tested.

\newpage

\hypertarget{superseded-phase-0m-r3a-summary}{%
\subsection{1.4 Superseded Phase-0M-R3A
summary}\label{superseded-phase-0m-r3a-summary}}

A narrow closure gate on four issues left open by Phase 0M-R3. Three
were confirmed as defects; one claim was found to have the wrong sign.

\begin{longtable}[]{@{}ll@{}}
\toprule
\begin{minipage}[b]{0.47\columnwidth}\raggedright
question\strut
\end{minipage} & \begin{minipage}[b]{0.47\columnwidth}\raggedright
answer\strut
\end{minipage}\tabularnewline
\midrule
\endhead
\begin{minipage}[t]{0.47\columnwidth}\raggedright
Was the 15 -\textgreater{} 20 um routing repair source-backed?\strut
\end{minipage} & \begin{minipage}[t]{0.47\columnwidth}\raggedright
\textbf{No --- DEFECT \#58.} The range quoted to license it,
\texttt{tsv.radius\_m}, governs THROUGH-SILICON vias. \texttt{ls\_to\_L}
is a \textbf{673 um} stacked via between embedded-component layers
inside the converter cell, where the Z1 TSV field spans 74 um. The
design space holds \texttt{tsv.*}, \texttt{rdl.*} and \texttt{bump.*}
and has no entry for this structure. The repair is a design point in an
unsourced process window\strut
\end{minipage}\tabularnewline
\begin{minipage}[t]{0.47\columnwidth}\raggedright
Is the repair arithmetic right?\strut
\end{minipage} & \begin{minipage}[t]{0.47\columnwidth}\raggedright
\textbf{Yes.} An independent implementation sharing no code with
production reproduces the routing ratio exactly at 5 of 5 radii ---
1.2360 at r = 15 um, 0.9344 at r = 20 um. The numbers were right; only
their provenance was not\strut
\end{minipage}\tabularnewline
\begin{minipage}[t]{0.47\columnwidth}\raggedright
Does a 34 um TSV cost 22.8 \% more vias?\strut
\end{minipage} & \begin{minipage}[t]{0.47\columnwidth}\raggedright
\textbf{No --- DEFECT \#59, the sign was backwards.} n\_required scales
as 1/A\_cu, so 30 -\textgreater{} 34 um gives -22.8 \%: 15140
-\textgreater{} 11692. It is a reduction, not a cost. The sweep also
skipped \textbf{32 um}, which passes the same bound at M\_AR =
1.067\strut
\end{minipage}\tabularnewline
\begin{minipage}[t]{0.47\columnwidth}\raggedright
Was the current-state table current?\strut
\end{minipage} & \begin{minipage}[t]{0.47\columnwidth}\raggedright
\textbf{No --- DEFECT \#57.} A section headed ``Current scientific
status (Phase 0M-R3)'' carried a table saying phase = 0M-R2 complete,
the 0M-R2 verdict, ZERO OF TWELVE REALIZABLE, and two already-completed
actions listed as future work\strut
\end{minipage}\tabularnewline
\bottomrule
\end{longtable}

Nothing was re-solved. The 30 -\textgreater{} 34 um change moves every
E/T/M property by under 0.1 \%, because under minimum-capability
provisioning the copper fraction is
\texttt{I\ /\ (J\_design\ *\ A\_field)}, in which the radius cancels.

\textbf{Verdict B --- the Z1/Z2 ordering survives, but the intra-VR via
technology remains unsourced.} Phase 0M-R3's verdict rested on package
clearances being the sole remaining blocker. DEFECT \#58 adds a second,
independent unsourced manufacturing assumption, so that premise no
longer holds.

\newpage

\hypertarget{superseded-phase-0m-r3-summary}{%
\subsection{1.5 Superseded Phase-0M-R3
summary}\label{superseded-phase-0m-r3-summary}}

\begin{longtable}[]{@{}ll@{}}
\toprule
\begin{minipage}[b]{0.47\columnwidth}\raggedright
question\strut
\end{minipage} & \begin{minipage}[b]{0.47\columnwidth}\raggedright
answer\strut
\end{minipage}\tabularnewline
\midrule
\endhead
\begin{minipage}[t]{0.47\columnwidth}\raggedright
Is there a buildable Z1 and a buildable Z2 in the existing
library?\strut
\end{minipage} & \begin{minipage}[t]{0.47\columnwidth}\raggedright
\textbf{Yes, both CONDITIONAL\_REALIZABLE}, passing every condition of a
contract frozen before the search, at all 27 corners of the bounded
clearance envelope. Neither is PROVEN\_REALIZABLE: three package
clearances are unsourced\strut
\end{minipage}\tabularnewline
\begin{minipage}[t]{0.47\columnwidth}\raggedright
What repair was needed?\strut
\end{minipage} & \begin{minipage}[t]{0.47\columnwidth}\raggedright
\textbf{One.} The intra-VR via radius 15 -\textgreater{} 20 um at the
framework's own legal pitch. \emph{(Phase 0M-R3A withdrew ``inside the
sourced range'': that range governs through-silicon vias, DEFECT \#58.)}
The binding net was \texttt{ls\_to\_L}, a vertical via stack at 1.236x
the vias available\strut
\end{minipage}\tabularnewline
\begin{minipage}[t]{0.47\columnwidth}\raggedright
Does the ordering survive?\strut
\end{minipage} & \begin{minipage}[t]{0.47\columnwidth}\raggedright
\textbf{Yes, essentially unchanged.} D\_driver 1.1942 -\textgreater{}
1.1943, +0.009 \%\strut
\end{minipage}\tabularnewline
\begin{minipage}[t]{0.47\columnwidth}\raggedright
Does the 30 um TSV have process margin?\strut
\end{minipage} & \begin{minipage}[t]{0.47\columnwidth}\raggedright
\textbf{No.} ARCH\_Z2 sits at M\_AR = 1.000 and fails a modest tolerance
bound. 34 um restores 1.089. \emph{(Phase 0M-R3A: that is 22.8 \% FEWER
vias, not more, and 32 um already passes at M\_AR = 1.067; DEFECT
\#59.)} Every E/T/M property moves by \textless{}= 0.02 \%\strut
\end{minipage}\tabularnewline
\begin{minipage}[t]{0.47\columnwidth}\raggedright
Was a transitional cooling correlation found?\strut
\end{minipage} & \begin{minipage}[t]{0.47\columnwidth}\raggedright
\textbf{No.} 2300 \textless{} Re \textless{} 3000 is permanently
OUT\_OF\_MODEL\_DOMAIN and the design space is capped at Re \textless{}=
2300\strut
\end{minipage}\tabularnewline
\begin{minipage}[t]{0.47\columnwidth}\raggedright
Was the previous zero-stale certificate right?\strut
\end{minipage} & \begin{minipage}[t]{0.47\columnwidth}\raggedright
\textbf{No.} 18 findings; DEFECT \#56\strut
\end{minipage}\tabularnewline
\bottomrule
\end{longtable}

\newpage

\hypertarget{superseded-phase-0m-r2-summary}{%
\subsection{1.6 Superseded Phase-0M-R2
summary}\label{superseded-phase-0m-r2-summary}}

This gate asked five falsification questions. Four returned a defect.

\begin{longtable}[]{@{}ll@{}}
\toprule
\begin{minipage}[b]{0.47\columnwidth}\raggedright
question\strut
\end{minipage} & \begin{minipage}[b]{0.47\columnwidth}\raggedright
answer\strut
\end{minipage}\tabularnewline
\midrule
\endhead
\begin{minipage}[t]{0.47\columnwidth}\raggedright
Was the previous phase's ``0 stale current-facing claims'' true?\strut
\end{minipage} & \begin{minipage}[t]{0.47\columnwidth}\raggedright
\textbf{No.} 17 stale findings across 8 claims, behind a green 46/46
cross-report run. DEFECT \#54: the checker matched each report as one
flat string and could not detect an omission at all\strut
\end{minipage}\tabularnewline
\begin{minipage}[t]{0.47\columnwidth}\raggedright
Is ARCH\_Z4 a real power die?\strut
\end{minipage} & \begin{minipage}[t]{0.47\columnwidth}\raggedright
\textbf{Now yes, still conditional.} 1921 um of sourced components
replace a heat source in a 10 um underfill layer. Its construction is
unsourced, and Z4 takes first or second place depending on which
construction is real\strut
\end{minipage}\tabularnewline
\begin{minipage}[t]{0.47\columnwidth}\raggedright
Does each TSV span what the architecture says?\strut
\end{minipage} & \begin{minipage}[t]{0.47\columnwidth}\raggedright
\textbf{Now yes.} Three domains had been using 300, 280 and a hard-coded
ladder for the same via. No architecture traverses 1300 um; ARCH\_Z1
traverses 74 um. Three of four sit exactly on the 10:1 limit\strut
\end{minipage}\tabularnewline
\begin{minipage}[t]{0.47\columnwidth}\raggedright
Is the cooler inside its correlation domain?\strut
\end{minipage} & \begin{minipage}[t]{0.47\columnwidth}\raggedright
\textbf{At 27 g/s yes, at 50 g/s no.} Re = 2342 lies between the laminar
ceiling of 2300 and Gnielinski's floor of 3000. h = 166 000 W/m2K, 1.46
MPa and 74 W are withdrawn\strut
\end{minipage}\tabularnewline
\begin{minipage}[t]{0.47\columnwidth}\raggedright
Does the RDL thermal tensor change the answer?\strut
\end{minipage} & \begin{minipage}[t]{0.47\columnwidth}\raggedright
\textbf{It changes absolute temperature by 7 K and the ordering by 0.105
\%.} Promoted to the production default for the absolute numbers\strut
\end{minipage}\tabularnewline
\bottomrule
\end{longtable}

A sixth finding came out of the second: \textbf{DEFECT \#55}, the
TSV-field material override was keyed on in-plane footprint, so the PCB,
the BGA field and ARCH\_Z2's backside converter were all modelled as a
copper/silicon composite -- k too high by 530x, E by 6.3x, CTE too low
by 3.4x, in the layers that set package warpage. It was
architecture-dependent and did not cancel: correcting it moved the Z1/Z2
driver ratio from 1.2552 to 1.1942.

\textbf{The deciding ordering survived all of it.} driver\_Z1 /
driver\_Z2 = 1.1942, ARCH\_Z2 better. \textbf{Nothing is buildable:} 0
of 12 template-architecture combinations are realizable at any of 27
clearance-bound corners, and area, routing and current density block
more of them than the clearances do.

\newpage

\hypertarget{the-physical-system-and-the-architectures}{%
\section{2. The physical system and the
architectures}\label{the-physical-system-and-the-architectures}}

\begin{longtable}[]{@{}lll@{}}
\toprule
\begin{minipage}[b]{0.30\columnwidth}\raggedright
quantity\strut
\end{minipage} & \begin{minipage}[b]{0.30\columnwidth}\raggedright
value\strut
\end{minipage} & \begin{minipage}[b]{0.30\columnwidth}\raggedright
provenance\strut
\end{minipage}\tabularnewline
\midrule
\endhead
\begin{minipage}[t]{0.30\columnwidth}\raggedright
load\strut
\end{minipage} & \begin{minipage}[t]{0.30\columnwidth}\raggedright
1000 W at 1 V, 1000 A\strut
\end{minipage} & \begin{minipage}[t]{0.30\columnwidth}\raggedright
reference system (June-2026)\strut
\end{minipage}\tabularnewline
\begin{minipage}[t]{0.30\columnwidth}\raggedright
input rail\strut
\end{minipage} & \begin{minipage}[t]{0.30\columnwidth}\raggedright
48 V\strut
\end{minipage} & \begin{minipage}[t]{0.30\columnwidth}\raggedright
reference system\strut
\end{minipage}\tabularnewline
\begin{minipage}[t]{0.30\columnwidth}\raggedright
compute die\strut
\end{minipage} & \begin{minipage}[t]{0.30\columnwidth}\raggedright
22.36 mm square, 500 mm2\strut
\end{minipage} & \begin{minipage}[t]{0.30\columnwidth}\raggedright
reference system\strut
\end{minipage}\tabularnewline
\begin{minipage}[t]{0.30\columnwidth}\raggedright
voltage regulators\strut
\end{minipage} & \begin{minipage}[t]{0.30\columnwidth}\raggedright
49 cells, 20.41 A each\strut
\end{minipage} & \begin{minipage}[t]{0.30\columnwidth}\raggedright
Phase-0F 1 kW design point\strut
\end{minipage}\tabularnewline
\begin{minipage}[t]{0.30\columnwidth}\raggedright
converter topology\strut
\end{minipage} & \begin{minipage}[t]{0.30\columnwidth}\raggedright
3:1 switched capacitor into a 3-phase buck (\texttt{N3K3KA2})\strut
\end{minipage} & \begin{minipage}[t]{0.30\columnwidth}\raggedright
compiled by the upstream optimizer\strut
\end{minipage}\tabularnewline
\begin{minipage}[t]{0.30\columnwidth}\raggedright
switching frequency\strut
\end{minipage} & \begin{minipage}[t]{0.30\columnwidth}\raggedright
805.17 kHz\strut
\end{minipage} & \begin{minipage}[t]{0.30\columnwidth}\raggedright
compiled\strut
\end{minipage}\tabularnewline
\begin{minipage}[t]{0.30\columnwidth}\raggedright
devices\strut
\end{minipage} & \begin{minipage}[t]{0.30\columnwidth}\raggedright
EPC2014C high side, EPC2057 low side\strut
\end{minipage} & \begin{minipage}[t]{0.30\columnwidth}\raggedright
EPC product table\strut
\end{minipage}\tabularnewline
\begin{minipage}[t]{0.30\columnwidth}\raggedright
inductor\strut
\end{minipage} & \begin{minipage}[t]{0.30\columnwidth}\raggedright
250 nH, 1.6 mOhm DCR, 500 um tall\strut
\end{minipage} & \begin{minipage}[t]{0.30\columnwidth}\raggedright
Rasheedi, TCPMT 2026\strut
\end{minipage}\tabularnewline
\begin{minipage}[t]{0.30\columnwidth}\raggedright
capacitor\strut
\end{minipage} & \begin{minipage}[t]{0.30\columnwidth}\raggedright
10 uF, 2.5 mOhm ESR\strut
\end{minipage} & \begin{minipage}[t]{0.30\columnwidth}\raggedright
Murata GCM21 datasheet\strut
\end{minipage}\tabularnewline
\bottomrule
\end{longtable}

Four vertical-integration architectures are compared. They differ in
\textbf{where the conversion happens}, which decides \textbf{which rail
crosses the substrate} --- and that single fact dominates the electrical
comparison.

\begin{longtable}[]{@{}llll@{}}
\toprule
\begin{minipage}[b]{0.22\columnwidth}\raggedright
architecture\strut
\end{minipage} & \begin{minipage}[b]{0.22\columnwidth}\raggedright
where the converter sits\strut
\end{minipage} & \begin{minipage}[b]{0.22\columnwidth}\raggedright
rail on the TSVs\strut
\end{minipage} & \begin{minipage}[b]{0.22\columnwidth}\raggedright
substrate\strut
\end{minipage}\tabularnewline
\midrule
\endhead
\begin{minipage}[t]{0.22\columnwidth}\raggedright
\textbf{Z1 EMBEDDED}\strut
\end{minipage} & \begin{minipage}[t]{0.22\columnwidth}\raggedright
inside the substrate power layers\strut
\end{minipage} & \begin{minipage}[t]{0.22\columnwidth}\raggedright
48 V, 23.4 A\strut
\end{minipage} & \begin{minipage}[t]{0.22\columnwidth}\raggedright
1.30 mm (component-height limited)\strut
\end{minipage}\tabularnewline
\begin{minipage}[t]{0.22\columnwidth}\raggedright
\textbf{Z2 BACKSIDE}\strut
\end{minipage} & \begin{minipage}[t]{0.22\columnwidth}\raggedright
on the substrate underside, as chiplets\strut
\end{minipage} & \begin{minipage}[t]{0.22\columnwidth}\raggedright
\textbf{1 V, 1000 A}\strut
\end{minipage} & \begin{minipage}[t]{0.22\columnwidth}\raggedright
0.30 mm\strut
\end{minipage}\tabularnewline
\begin{minipage}[t]{0.22\columnwidth}\raggedright
\textbf{Z4 POWER DIE}\strut
\end{minipage} & \begin{minipage}[t]{0.22\columnwidth}\raggedright
a dedicated die between substrate and compute die\strut
\end{minipage} & \begin{minipage}[t]{0.22\columnwidth}\raggedright
48 V, 23.4 A\strut
\end{minipage} & \begin{minipage}[t]{0.22\columnwidth}\raggedright
0.30 mm\strut
\end{minipage}\tabularnewline
\begin{minipage}[t]{0.22\columnwidth}\raggedright
\textbf{Z5 2.5D LATERAL}\strut
\end{minipage} & \begin{minipage}[t]{0.22\columnwidth}\raggedright
beside the die on a common interposer\strut
\end{minipage} & \begin{minipage}[t]{0.22\columnwidth}\raggedright
48 V\strut
\end{minipage} & \begin{minipage}[t]{0.22\columnwidth}\raggedright
0.30 mm\strut
\end{minipage}\tabularnewline
\bottomrule
\end{longtable}

Z2 puts 1000 A on the TSV field where Z1 and Z4 put 23.4 A: a factor of
\textbf{1822 in \(I^2\)} at the same geometry. Z5 is the control --- it
is the architecture DVPD exists to beat.

\begin{figure}
\centering
\includegraphics{/home/beams0/MGHARIB/Mohamed_stuff/HVR_OPT/results/phase0m/FIG_0M_E_partition_templates.png}
\caption{Architecture cross-sections}
\end{figure}

\emph{Component-partition templates and whether each fits the die
shadow.}

\newpage

\hypertarget{the-data-flow-end-to-end}{%
\section{3. The data flow, end to end}\label{the-data-flow-end-to-end}}

\begin{verbatim}
  architecture choice
        |
        v
  component partition template          which component sits on which tier
        |
        v
  converter synthesis (upstream)        devices, inductor, capacitor, f_sw, area
        |
        v
  loss ledger                           one owner per watt, system scope enforced
        |
        +--> intra-VR routing           metal BETWEEN the converter components
        +--> vertical interconnect      R_bulk + dR_transition, per class
        +--> package RDL                multilayer resistive sheet
        |
        v
  spatial heat deposition               each term into the layer it occurs in
        |                               (checked against the electrical total)
        v
  thermal solve (CalculiX)  <---->  physical cooling
        |                            h(x) from channel geometry and flow
        |                            T_fluid(x) from the heat removed upstream
        v
  temperature field T(x,y,z)
        |
        +--> back to the loss ledger    rho(T), R_DS(on)(T), DCR(T)   [FIXED POINT]
        |
        v
  structural solve (CalculiX)           global thermoelastic warpage and stress
        |
        v
  corner-joint relative displacement    the driver handed to the local model
        |
        v
  local solder submodel (CalculiX)      viscoplastic, cyclic, graded mesh
        |
        v
  per-cycle local inelastic RESPONSE metric --> RELATIVE life only
\end{verbatim}

Two loops close and one does not. \textbf{Electrical and thermal are
iterated to a fixed point.} \textbf{Thermal drives mechanics one way.}
Mechanics returns nothing, and Section 19 gives the measured reason.

\begin{figure}
\centering
\includegraphics{/home/beams0/MGHARIB/Mohamed_stuff/HVR_OPT/results/phase0mr/FIG_0MR_A_production_flow.png}
\caption{Production flow}
\end{figure}

\emph{Which parts of this flow were measured in Phase 0M and which now
drive it.}

\newpage

\hypertarget{governing-equations}{%
\section{4. Governing equations}\label{governing-equations}}

\hypertarget{electrical}{%
\subsection{4.1 Electrical}\label{electrical}}

Conductor resistance, temperature dependent:

\[R = \frac{\rho(T)\,\ell}{A}, \qquad \rho(T) = \rho_{20}\left[1 + \alpha (T - 20)\right]\]

Via-field allocation and field resistance, with the count DISCRETE:

\[n_{req} = \left\lceil \frac{I}{J_{design} A_c} \right\rceil, \qquad R_{field} = \frac{R_{single}}{n_{alloc}} + \Delta R_{transition}\]

Under minimum-capability provisioning the first term collapses to
\(R_{field} = \rho(T) h J_{design} / I\) and \textbf{the conductor
cross-section cancels}. That is derived, verified numerically and
discussed in Section 9.

Lateral distribution, solved as a PDE rather than a lumped resistance:

\[\nabla \cdot \left( \sigma_s \nabla V \right) = -J, \qquad \sigma_s = \frac{N_{layers}\, t_{Cu}\, \text{occupancy}}{\rho(T)}\]

Three-dimensional conduction, for transitions:

\[\nabla \cdot \left( \sigma \nabla V \right) = 0, \qquad \mathbf{J} = -\sigma \nabla V, \qquad P = \int \mathbf{J} \cdot \mathbf{E}\, dV\]

Converter loss --- the upstream closed form, decomposed per component
and checked against its own aggregate:

\[P_{cond} = \frac{V_{out}}{V_{in}} I_{rms}^2 \left[ R_{DS,hs} + 2\left(1-\tfrac{1}{N}\right) R_{C,ESR} + \left( \frac{V_{in}}{N^2 V_{out}}(R_{L} + R_{DS,ls}) + \frac{2 R_{DS,ls}}{N} \right) \Sigma B^2 - \text{corr} \right]\]

\hypertarget{thermal}{%
\subsection{4.2 Thermal}\label{thermal}}

\[\nabla \cdot \left( \mathbf{k}_{eff} \nabla T \right) + \dot{q} = 0\]

with an anisotropic \(\mathbf{k}_{eff}\) for every heterogeneous layer,
and the cooled boundary

\[-\mathbf{k}\nabla T \cdot \mathbf{n} = h(x)\left[T_s(x) - T_f(x)\right], \qquad T_f(x) = T_{in} + \frac{Q(0 \to x)}{\dot{m} c_p}\]

\(h\) from a Nusselt correlation, \(T_f\) from the heat actually removed
upstream. Both are iterated with the wall temperature to a fixed point.

\hypertarget{mechanical}{%
\subsection{4.3 Mechanical}\label{mechanical}}

Global, linear thermoelastic:

\[\nabla \cdot \boldsymbol{\sigma} = 0, \qquad \boldsymbol{\sigma} = \mathbf{C} : \left( \boldsymbol{\varepsilon} - \boldsymbol{\alpha}\,\Delta T \right)\]

Local solder, viscoplastic. \textbf{What CalculiX actually integrates in
a \texttt{*VISCO} step is an OVERSTRESS law}, and the distinction
matters:

\[\dot{\varepsilon}_{in} = A(T)\left( \sigma_{vm} - \sigma_y(\bar{\varepsilon}_{in}) \right)^{n(T)} \ \text{for}\ \sigma_{vm} > \sigma_y, \qquad \dot{\varepsilon}_{in} = 0 \ \text{otherwise}\]

There is \textbf{no rate-independent plastic strain} in such a step, and
\(A\), \(n\) and \(\sigma_y\) are interpolated \textbf{linearly} between
card rows by the solver --- not by an Arrhenius law. Any Arrhenius
relation exists only in how the rows were generated.

A separate rate-independent plastic branch \textbf{is not presently
justified if a validated UNIFIED viscoplastic law such as Anand is
used}, because a unified law already carries the rate-independent
response inside its single inelastic flow rule. That is a statement
about unified laws, not a general claim about plasticity.

\textbf{The variable mismatch, stated plainly.} The creep card in this
repository was fitted against \textbf{total applied stress}
\(\sigma_{vm}\) in the source experiments, while the solver drives it
with the \textbf{overstress} \(\sigma_{vm} - \sigma_y\). The two agree
only where \(\sigma_y\) is small against \(\sigma_{vm}\). This is not a
coding error and cannot be fixed by refitting the same data; it is why
the constitutive status is CONDITIONAL rather than closed, and it is the
dominant open blocker on every absolute reliability number in this
framework.

Damage metric, per stabilised cycle:

\[D = \Delta \text{PEEQ}\big|_{cycle}, \qquad W_{in} = \oint \boldsymbol{\sigma} : d\boldsymbol{\varepsilon}_{in}, \qquad \boldsymbol{\varepsilon}_{in} = \mathbf{ME} - \mathbf{C}^{-1}(T):\boldsymbol{\sigma}\]

\texttt{ME} is the solver's own mechanical strain. Reconstructing the
inelastic strain from it, rather than from a separately parsed
temperature, removes an entire class of error (Section 22, defect 46).

\newpage

\hypertarget{solvers-and-what-each-is-used-for}{%
\section{5. Solvers, and what each is used
for}\label{solvers-and-what-each-is-used-for}}

\begin{longtable}[]{@{}llll@{}}
\toprule
\begin{minipage}[b]{0.22\columnwidth}\raggedright
solver\strut
\end{minipage} & \begin{minipage}[b]{0.22\columnwidth}\raggedright
version\strut
\end{minipage} & \begin{minipage}[b]{0.22\columnwidth}\raggedright
used for\strut
\end{minipage} & \begin{minipage}[b]{0.22\columnwidth}\raggedright
why this one\strut
\end{minipage}\tabularnewline
\midrule
\endhead
\begin{minipage}[t]{0.22\columnwidth}\raggedright
\textbf{CalculiX}\strut
\end{minipage} & \begin{minipage}[t]{0.22\columnwidth}\raggedright
2.23\strut
\end{minipage} & \begin{minipage}[t]{0.22\columnwidth}\raggedright
thermal conduction, global thermoelasticity, local nonlinear
solder\strut
\end{minipage} & \begin{minipage}[t]{0.22\columnwidth}\raggedright
the primary physics solver. Practical ceiling about 40 000 elements for
this nonlinear transient class\strut
\end{minipage}\tabularnewline
\begin{minipage}[t]{0.22\columnwidth}\raggedright
\textbf{MOOSE} \texttt{heat\_transfer}\strut
\end{minipage} & \begin{minipage}[t]{0.22\columnwidth}\raggedright
09515f92aa\strut
\end{minipage} & \begin{minipage}[t]{0.22\columnwidth}\raggedright
independent cross-check on scalar conduction\strut
\end{minipage} & \begin{minipage}[t]{0.22\columnwidth}\raggedright
a different mesher, element formulation, quadrature and linear solver.
\texttt{div(k\ grad\ T)} and \texttt{div(sigma\ grad\ V)} are the same
operator\strut
\end{minipage}\tabularnewline
\begin{minipage}[t]{0.22\columnwidth}\raggedright
\textbf{finite-volume solver}\strut
\end{minipage} & \begin{minipage}[t]{0.22\columnwidth}\raggedright
this repository\strut
\end{minipage} & \begin{minipage}[t]{0.22\columnwidth}\raggedright
3-D DC conduction for transitions\strut
\end{minipage} & \begin{minipage}[t]{0.22\columnwidth}\raggedright
no DC conduction solver is installed. Cell-centred FV with harmonic-mean
face conductivities is exactly conservative across a nine-decade
conductivity jump\strut
\end{minipage}\tabularnewline
\begin{minipage}[t]{0.22\columnwidth}\raggedright
\textbf{analytic Python}\strut
\end{minipage} & \begin{minipage}[t]{0.22\columnwidth}\raggedright
--\strut
\end{minipage} & \begin{minipage}[t]{0.22\columnwidth}\raggedright
via allocation, sheet conduction, effective media, routing, cooling
correlations\strut
\end{minipage} & \begin{minipage}[t]{0.22\columnwidth}\raggedright
cheap enough to sweep, and each form is checked against a numerical
solve\strut
\end{minipage}\tabularnewline
\begin{minipage}[t]{0.22\columnwidth}\raggedright
\textbf{HFSS}\strut
\end{minipage} & \begin{minipage}[t]{0.22\columnwidth}\raggedright
AEDT 2026 R1\strut
\end{minipage} & \begin{minipage}[t]{0.22\columnwidth}\raggedright
AC / skin validation only\strut
\end{minipage} & \begin{minipage}[t]{0.22\columnwidth}\raggedright
licensed and available. It is a full-wave solver and \textbf{cannot pose
a pure conduction problem with no return path}, so it is NOT used for
the DC transition resistance\strut
\end{minipage}\tabularnewline
\begin{minipage}[t]{0.22\columnwidth}\raggedright
\textbf{Icepak}\strut
\end{minipage} & \begin{minipage}[t]{0.22\columnwidth}\raggedright
AEDT 2026 R1\strut
\end{minipage} & \begin{minipage}[t]{0.22\columnwidth}\raggedright
nothing, currently\strut
\end{minipage} & \begin{minipage}[t]{0.22\columnwidth}\raggedright
licensed and available. An earlier attempt returned a solver clamp of
exactly 5000 K on every object and that result stays discarded.
\textbf{No Icepak validation is claimed}\strut
\end{minipage}\tabularnewline
\begin{minipage}[t]{0.22\columnwidth}\raggedright
\textbf{Q3D}\strut
\end{minipage} & \begin{minipage}[t]{0.22\columnwidth}\raggedright
--\strut
\end{minipage} & \begin{minipage}[t]{0.22\columnwidth}\raggedright
nothing\strut
\end{minipage} & \begin{minipage}[t]{0.22\columnwidth}\raggedright
\textbf{not licensed}. Checked against the licence manager, not
assumed\strut
\end{minipage}\tabularnewline
\begin{minipage}[t]{0.22\columnwidth}\raggedright
Elmer, GetDP, Palace, OpenFOAM\strut
\end{minipage} & \begin{minipage}[t]{0.22\columnwidth}\raggedright
--\strut
\end{minipage} & \begin{minipage}[t]{0.22\columnwidth}\raggedright
nothing\strut
\end{minipage} & \begin{minipage}[t]{0.22\columnwidth}\raggedright
\textbf{not installed}. Palace's electrostatic mode solves a
permittivity problem and would not be a DC conduction solver in any
case\strut
\end{minipage}\tabularnewline
\bottomrule
\end{longtable}

\hypertarget{what-the-finite-volume-solver-was-verified-against}{%
\subsection{5.1 What the finite-volume solver was verified
against}\label{what-the-finite-volume-solver-was-verified-against}}

\begin{longtable}[]{@{}ll@{}}
\toprule
\begin{minipage}[b]{0.47\columnwidth}\raggedright
case\strut
\end{minipage} & \begin{minipage}[b]{0.47\columnwidth}\raggedright
result\strut
\end{minipage}\tabularnewline
\midrule
\endhead
\begin{minipage}[t]{0.47\columnwidth}\raggedright
uniform bar, analytic rho L / A\strut
\end{minipage} & \begin{minipage}[t]{0.47\columnwidth}\raggedright
\textbf{exact at every mesh}\strut
\end{minipage}\tabularnewline
\begin{minipage}[t]{0.47\columnwidth}\raggedright
circular contact on a half-space, analytic rho/(4a)\strut
\end{minipage} & \begin{minipage}[t]{0.47\columnwidth}\raggedright
+0.9 \% at 6 cells per contact radius\strut
\end{minipage}\tabularnewline
\begin{minipage}[t]{0.47\columnwidth}\raggedright
five transition geometries against MOOSE\strut
\end{minipage} & \begin{minipage}[t]{0.47\columnwidth}\raggedright
\textbf{0.25 to 0.73 \%, 5 of 5 ACCEPTED}\strut
\end{minipage}\tabularnewline
\bottomrule
\end{longtable}

Both terminals are equipotential Dirichlet planes in both codes, the
dielectric is present in both, and each run reports its discretised
block volumes so the geometry match is proved rather than assumed.

\newpage

\hypertarget{how-to-read-every-number-in-this-framework}{%
\section{6. How to read every number in this
framework}\label{how-to-read-every-number-in-this-framework}}

Each quantity carries one of these classes, and they are not
interchangeable:

\begin{longtable}[]{@{}ll@{}}
\toprule
\begin{minipage}[b]{0.47\columnwidth}\raggedright
class\strut
\end{minipage} & \begin{minipage}[b]{0.47\columnwidth}\raggedright
meaning\strut
\end{minipage}\tabularnewline
\midrule
\endhead
\begin{minipage}[t]{0.47\columnwidth}\raggedright
\textbf{SOURCE\_BACKED}\strut
\end{minipage} & \begin{minipage}[t]{0.47\columnwidth}\raggedright
taken from a named primary source with its conditions\strut
\end{minipage}\tabularnewline
\begin{minipage}[t]{0.47\columnwidth}\raggedright
\textbf{ANALYTICALLY\_DERIVED}\strut
\end{minipage} & \begin{minipage}[t]{0.47\columnwidth}\raggedright
a closed form, verified numerically\strut
\end{minipage}\tabularnewline
\begin{minipage}[t]{0.47\columnwidth}\raggedright
\textbf{NUMERICALLY\_VALIDATED}\strut
\end{minipage} & \begin{minipage}[t]{0.47\columnwidth}\raggedright
solved, mesh-converged, checked against a closed form\strut
\end{minipage}\tabularnewline
\begin{minipage}[t]{0.47\columnwidth}\raggedright
\textbf{INDEPENDENTLY\_CROSS\_VALIDATED}\strut
\end{minipage} & \begin{minipage}[t]{0.47\columnwidth}\raggedright
reproduced by a second, independent solver\strut
\end{minipage}\tabularnewline
\begin{minipage}[t]{0.47\columnwidth}\raggedright
\textbf{CALIBRATED}\strut
\end{minipage} & \begin{minipage}[t]{0.47\columnwidth}\raggedright
fitted to reproduce a target. Not transferable\strut
\end{minipage}\tabularnewline
\begin{minipage}[t]{0.47\columnwidth}\raggedright
\textbf{HOMOGENIZED}\strut
\end{minipage} & \begin{minipage}[t]{0.47\columnwidth}\raggedright
an effective property of a heterogeneous layer, with its RVE error\strut
\end{minipage}\tabularnewline
\begin{minipage}[t]{0.47\columnwidth}\raggedright
\textbf{CONDITIONAL}\strut
\end{minipage} & \begin{minipage}[t]{0.47\columnwidth}\raggedright
true only under a stated open input, and the input is named\strut
\end{minipage}\tabularnewline
\begin{minipage}[t]{0.47\columnwidth}\raggedright
\textbf{OPEN}\strut
\end{minipage} & \begin{minipage}[t]{0.47\columnwidth}\raggedright
not determined. Never silently defaulted\strut
\end{minipage}\tabularnewline
\bottomrule
\end{longtable}

The full ledger is \texttt{results/phase0mr/uncertainty\_ledger.csv}.
Its headline:

\begin{longtable}[]{@{}ll@{}}
\toprule
class & count\tabularnewline
\midrule
\endhead
INDEPENDENTLY\_CROSS\_VALIDATED & 1\tabularnewline
NUMERICALLY\_VALIDATED & 4\tabularnewline
ANALYTICALLY\_DERIVED & 2\tabularnewline
SOURCE\_BACKED & 1\tabularnewline
CONDITIONAL & 1\tabularnewline
\textbf{OPEN} & \textbf{5}\tabularnewline
\bottomrule
\end{longtable}

\textbf{Uncertainties are never RMS-summed across common-mode terms.}
For an architecture RATIO the cancellation is TESTED and the tested
residual is what is quoted.

\newpage

\hypertarget{input-parameters-and-their-provenance}{%
\section{7. Input parameters and their
provenance}\label{input-parameters-and-their-provenance}}

\begin{longtable}[]{@{}llll@{}}
\toprule
\begin{minipage}[b]{0.22\columnwidth}\raggedright
parameter\strut
\end{minipage} & \begin{minipage}[b]{0.22\columnwidth}\raggedright
value\strut
\end{minipage} & \begin{minipage}[b]{0.22\columnwidth}\raggedright
class\strut
\end{minipage} & \begin{minipage}[b]{0.22\columnwidth}\raggedright
source\strut
\end{minipage}\tabularnewline
\midrule
\endhead
\begin{minipage}[t]{0.22\columnwidth}\raggedright
rho\_Cu at 20 C\strut
\end{minipage} & \begin{minipage}[t]{0.22\columnwidth}\raggedright
17.0 nOhm.m\strut
\end{minipage} & \begin{minipage}[t]{0.22\columnwidth}\raggedright
SOURCE\_BACKED\strut
\end{minipage} & \begin{minipage}[t]{0.22\columnwidth}\raggedright
June-2026 Table III\strut
\end{minipage}\tabularnewline
\begin{minipage}[t]{0.22\columnwidth}\raggedright
rho\_solder at 20 C\strut
\end{minipage} & \begin{minipage}[t]{0.22\columnwidth}\raggedright
130 nOhm.m\strut
\end{minipage} & \begin{minipage}[t]{0.22\columnwidth}\raggedright
SOURCE\_BACKED\strut
\end{minipage} & \begin{minipage}[t]{0.22\columnwidth}\raggedright
June-2026\strut
\end{minipage}\tabularnewline
\begin{minipage}[t]{0.22\columnwidth}\raggedright
alpha\_Cu\strut
\end{minipage} & \begin{minipage}[t]{0.22\columnwidth}\raggedright
3.93e-3 /K\strut
\end{minipage} & \begin{minipage}[t]{0.22\columnwidth}\raggedright
ASSUMPTION\strut
\end{minipage} & \begin{minipage}[t]{0.22\columnwidth}\raggedright
bulk Cu TCR\strut
\end{minipage}\tabularnewline
\begin{minipage}[t]{0.22\columnwidth}\raggedright
TSV diameter / pitch / height\strut
\end{minipage} & \begin{minipage}[t]{0.22\columnwidth}\raggedright
30 / 60 / 300 um\strut
\end{minipage} & \begin{minipage}[t]{0.22\columnwidth}\raggedright
SOURCE\_BACKED\strut
\end{minipage} & \begin{minipage}[t]{0.22\columnwidth}\raggedright
June-2026 Table III\strut
\end{minipage}\tabularnewline
\begin{minipage}[t]{0.22\columnwidth}\raggedright
TSV liner thickness\strut
\end{minipage} & \begin{minipage}[t]{0.22\columnwidth}\raggedright
0.5 um\strut
\end{minipage} & \begin{minipage}[t]{0.22\columnwidth}\raggedright
ASSUMPTION\strut
\end{minipage} & \begin{minipage}[t]{0.22\columnwidth}\raggedright
typical SiO2 liner; swept\strut
\end{minipage}\tabularnewline
\begin{minipage}[t]{0.22\columnwidth}\raggedright
TSV aspect-ratio limit\strut
\end{minipage} & \begin{minipage}[t]{0.22\columnwidth}\raggedright
10:1\strut
\end{minipage} & \begin{minipage}[t]{0.22\columnwidth}\raggedright
DERIVED\_FROM\_SOURCE\strut
\end{minipage} & \begin{minipage}[t]{0.22\columnwidth}\raggedright
the source demonstrates d = 30 um at h = 300 um\strut
\end{minipage}\tabularnewline
\begin{minipage}[t]{0.22\columnwidth}\raggedright
keep-out ratio\strut
\end{minipage} & \begin{minipage}[t]{0.22\columnwidth}\raggedright
swept 0 to 1.5\strut
\end{minipage} & \begin{minipage}[t]{0.22\columnwidth}\raggedright
ASSUMPTION\strut
\end{minipage} & \begin{minipage}[t]{0.22\columnwidth}\raggedright
no source; never fixed\strut
\end{minipage}\tabularnewline
\begin{minipage}[t]{0.22\columnwidth}\raggedright
RDL stack\strut
\end{minipage} & \begin{minipage}[t]{0.22\columnwidth}\raggedright
t = 10 um, W/S = 10/2, 2 VDD + 2 GND\strut
\end{minipage} & \begin{minipage}[t]{0.22\columnwidth}\raggedright
SOURCE\_BACKED\strut
\end{minipage} & \begin{minipage}[t]{0.22\columnwidth}\raggedright
ECTC-2024 Table I\strut
\end{minipage}\tabularnewline
\begin{minipage}[t]{0.22\columnwidth}\raggedright
J\_design\strut
\end{minipage} & \begin{minipage}[t]{0.22\columnwidth}\raggedright
1.0e8 A/m2\strut
\end{minipage} & \begin{minipage}[t]{0.22\columnwidth}\raggedright
DERIVED\_FROM\_SOURCE\strut
\end{minipage} & \begin{minipage}[t]{0.22\columnwidth}\raggedright
June-2026 provisioning\strut
\end{minipage}\tabularnewline
\begin{minipage}[t]{0.22\columnwidth}\raggedright
J\_process\_max\strut
\end{minipage} & \begin{minipage}[t]{0.22\columnwidth}\raggedright
4.0e8 A/m2\strut
\end{minipage} & \begin{minipage}[t]{0.22\columnwidth}\raggedright
DERIVED\_FROM\_SOURCE\strut
\end{minipage} & \begin{minipage}[t]{0.22\columnwidth}\raggedright
ECTC-2024, the most aggressive available\strut
\end{minipage}\tabularnewline
\begin{minipage}[t]{0.22\columnwidth}\raggedright
\textbf{J\_reliability\_allow}\strut
\end{minipage} & \begin{minipage}[t]{0.22\columnwidth}\raggedright
--\strut
\end{minipage} & \begin{minipage}[t]{0.22\columnwidth}\raggedright
\textbf{OPEN}\strut
\end{minipage} & \begin{minipage}[t]{0.22\columnwidth}\raggedright
no primary EM study with its T, lifetime and geometry was
retrievable\strut
\end{minipage}\tabularnewline
\begin{minipage}[t]{0.22\columnwidth}\raggedright
GaN heights\strut
\end{minipage} & \begin{minipage}[t]{0.22\columnwidth}\raggedright
685 / 606 um\strut
\end{minipage} & \begin{minipage}[t]{0.22\columnwidth}\raggedright
SOURCE\_BACKED\strut
\end{minipage} & \begin{minipage}[t]{0.22\columnwidth}\raggedright
EPC2212 / EPC2090 datasheets\strut
\end{minipage}\tabularnewline
\begin{minipage}[t]{0.22\columnwidth}\raggedright
inductor\strut
\end{minipage} & \begin{minipage}[t]{0.22\columnwidth}\raggedright
250 nH, 1.6 mOhm, 500 um\strut
\end{minipage} & \begin{minipage}[t]{0.22\columnwidth}\raggedright
SOURCE\_BACKED\strut
\end{minipage} & \begin{minipage}[t]{0.22\columnwidth}\raggedright
Rasheedi, TCPMT 2026\strut
\end{minipage}\tabularnewline
\begin{minipage}[t]{0.22\columnwidth}\raggedright
MLCC height\strut
\end{minipage} & \begin{minipage}[t]{0.22\columnwidth}\raggedright
100 um nominal, 250 um upper\strut
\end{minipage} & \begin{minipage}[t]{0.22\columnwidth}\raggedright
PARTIAL\strut
\end{minipage} & \begin{minipage}[t]{0.22\columnwidth}\raggedright
manufacturer page unreachable; envelope carried\strut
\end{minipage}\tabularnewline
\begin{minipage}[t]{0.22\columnwidth}\raggedright
solder E, nu, CTE\strut
\end{minipage} & \begin{minipage}[t]{0.22\columnwidth}\raggedright
45 GPa, 0.36, 21.5 ppm/K\strut
\end{minipage} & \begin{minipage}[t]{0.22\columnwidth}\raggedright
ASSUMPTION\strut
\end{minipage} & \begin{minipage}[t]{0.22\columnwidth}\raggedright
commonly quoted band\strut
\end{minipage}\tabularnewline
\begin{minipage}[t]{0.22\columnwidth}\raggedright
\textbf{solder creep A(T), n(T)}\strut
\end{minipage} & \begin{minipage}[t]{0.22\columnwidth}\raggedright
19 points over 7.51-15.92 MPa\strut
\end{minipage} & \begin{minipage}[t]{0.22\columnwidth}\raggedright
\textbf{CONDITIONAL}\strut
\end{minipage} & \begin{minipage}[t]{0.22\columnwidth}\raggedright
\textbf{Sn3.9Ag0.6Cu, NOT SAC305}\strut
\end{minipage}\tabularnewline
\begin{minipage}[t]{0.22\columnwidth}\raggedright
underfill E / CTE\strut
\end{minipage} & \begin{minipage}[t]{0.22\columnwidth}\raggedright
8 GPa / 30 ppm/K\strut
\end{minipage} & \begin{minipage}[t]{0.22\columnwidth}\raggedright
PROVISIONAL\strut
\end{minipage} & \begin{minipage}[t]{0.22\columnwidth}\raggedright
unsourced; sensitivity run\strut
\end{minipage}\tabularnewline
\begin{minipage}[t]{0.22\columnwidth}\raggedright
board\strut
\end{minipage} & \begin{minipage}[t]{0.22\columnwidth}\raggedright
20 GPa, 17 ppm/K, 1.6 mm FR4\strut
\end{minipage} & \begin{minipage}[t]{0.22\columnwidth}\raggedright
ASSUMPTION\strut
\end{minipage} & \begin{minipage}[t]{0.22\columnwidth}\raggedright
typical laminate\strut
\end{minipage}\tabularnewline
\begin{minipage}[t]{0.22\columnwidth}\raggedright
coolant\strut
\end{minipage} & \begin{minipage}[t]{0.22\columnwidth}\raggedright
water, 27 g/s, 25 C inlet\strut
\end{minipage} & \begin{minipage}[t]{0.22\columnwidth}\raggedright
DESIGN VARIABLE\strut
\end{minipage} & \begin{minipage}[t]{0.22\columnwidth}\raggedright
derived, not calibrated\strut
\end{minipage}\tabularnewline
\begin{minipage}[t]{0.22\columnwidth}\raggedright
\textbf{required life}\strut
\end{minipage} & \begin{minipage}[t]{0.22\columnwidth}\raggedright
--\strut
\end{minipage} & \begin{minipage}[t]{0.22\columnwidth}\raggedright
\textbf{OPEN}\strut
\end{minipage} & \begin{minipage}[t]{0.22\columnwidth}\raggedright
no field or product requirement exists\strut
\end{minipage}\tabularnewline
\begin{minipage}[t]{0.22\columnwidth}\raggedright
\textbf{mechanical allowable}\strut
\end{minipage} & \begin{minipage}[t]{0.22\columnwidth}\raggedright
--\strut
\end{minipage} & \begin{minipage}[t]{0.22\columnwidth}\raggedright
\textbf{OPEN}\strut
\end{minipage} & \begin{minipage}[t]{0.22\columnwidth}\raggedright
nine candidates tested, none admissible\strut
\end{minipage}\tabularnewline
\bottomrule
\end{longtable}

\newpage

\hypertarget{interconnect-geometry}{%
\section{8. Interconnect geometry}\label{interconnect-geometry}}

\textbf{No resistance is a primitive.} Every interconnect quantity is
classified as exactly one of MATERIAL, PROCESS, DESIGN or DERIVED, and
the classifier raises on anything unclassified.

\begin{longtable}[]{@{}ll@{}}
\toprule
\begin{minipage}[b]{0.47\columnwidth}\raggedright
class\strut
\end{minipage} & \begin{minipage}[b]{0.47\columnwidth}\raggedright
examples\strut
\end{minipage}\tabularnewline
\midrule
\endhead
\begin{minipage}[t]{0.47\columnwidth}\raggedright
MATERIAL\strut
\end{minipage} & \begin{minipage}[t]{0.47\columnwidth}\raggedright
rho\_Cu(T), alpha, k\_Cu, k\_Si, k\_SiO2\strut
\end{minipage}\tabularnewline
\begin{minipage}[t]{0.47\columnwidth}\raggedright
PROCESS\strut
\end{minipage} & \begin{minipage}[t]{0.47\columnwidth}\raggedright
minimum pitch, liner thickness, aspect-ratio limit, keep-out ratio,
J\_process\_max\strut
\end{minipage}\tabularnewline
\begin{minipage}[t]{0.47\columnwidth}\raggedright
DESIGN\strut
\end{minipage} & \begin{minipage}[t]{0.47\columnwidth}\raggedright
via radius, pitch, height, allocated count, RDL thickness / width /
spacing / layer count, bump diameter / height / pitch, pad size\strut
\end{minipage}\tabularnewline
\begin{minipage}[t]{0.47\columnwidth}\raggedright
DERIVED\strut
\end{minipage} & \begin{minipage}[t]{0.47\columnwidth}\raggedright
conductor area, single-via R, per-via current limit, available and
required counts, field R, Joule loss, copper fraction, effective k and
CTE, transition dR\strut
\end{minipage}\tabularnewline
\bottomrule
\end{longtable}

Three quantities that used to be primitives and are now derived:

\begin{enumerate}
\def\labelenumi{\arabic{enumi}.}
\tightlist
\item
  the per-via current limit, which was entered in amperes and is now
  \texttt{J\_design\ *\ A\_c};
\item
  the current-crowding correction, which was one global constant and is
  now solved per class and per geometry;
\item
  the TSV liner, which was absent and now sets the conducting radius.
\end{enumerate}

\hypertarget{transition-resistance}{%
\subsection{8.1 Transition resistance}\label{transition-resistance}}

\[R_{total} = R_{bulk}(\text{geometry}, T) + \Delta R_{transition}(\text{class}, \text{geometry}, T)\]

\texttt{dR} is solved exactly for each production geometry by the
finite-volume solver and cached against a geometry fingerprint. There is
\textbf{no fitted surrogate}: a two-term logarithmic form was tried over
178 geometries and rejected at 9 to 29 \% of the spread of the quantity
it was fitting.

\textbf{\texttt{dR} can be negative, and that is physical, not an
error.} The 1-D reference is a series of named conductors with implicit
equipotential cross-sections everywhere; the 3-D solve constrains only
the two terminal planes. A generous landing pad genuinely replaces part
of a narrow trace. \texttt{dR} is therefore reported as a difference
against a stated reference and never as an ``excess''.

\begin{longtable}[]{@{}llll@{}}
\toprule
class & R\_1D & dR & R\_3D / R\_1D\tabularnewline
\midrule
\endhead
bga & 2.5253 mOhm & -0.3255 mOhm & \textbf{0.8711}\tabularnewline
bga & 2.5253 mOhm & -0.3255 mOhm & \textbf{0.8711}\tabularnewline
bga & 2.5253 mOhm & -0.3255 mOhm & \textbf{0.8711}\tabularnewline
bga & 2.5253 mOhm & -0.3255 mOhm & \textbf{0.8711}\tabularnewline
bga & 2.5253 mOhm & -0.3255 mOhm & \textbf{0.8711}\tabularnewline
bga & 2.5253 mOhm & -0.3255 mOhm & \textbf{0.8711}\tabularnewline
bga & 2.5253 mOhm & -0.3255 mOhm & \textbf{0.8711}\tabularnewline
bga & 2.5253 mOhm & -0.3255 mOhm & \textbf{0.8711}\tabularnewline
bga & 2.5253 mOhm & -0.3255 mOhm & \textbf{0.8711}\tabularnewline
bga & 2.5253 mOhm & -0.3255 mOhm & \textbf{0.8711}\tabularnewline
bga & 2.5253 mOhm & -0.3255 mOhm & \textbf{0.8711}\tabularnewline
bga & 2.5253 mOhm & -0.3255 mOhm & \textbf{0.8711}\tabularnewline
bga & 2.5253 mOhm & -0.3255 mOhm & \textbf{0.8711}\tabularnewline
bga & 2.5253 mOhm & -0.3255 mOhm & \textbf{0.8711}\tabularnewline
bga & 2.5253 mOhm & -0.3255 mOhm & \textbf{0.8711}\tabularnewline
bga & 2.5253 mOhm & -0.3255 mOhm & \textbf{0.8711}\tabularnewline
bga & 2.5253 mOhm & -0.3255 mOhm & \textbf{0.8711}\tabularnewline
bga & 2.5253 mOhm & -0.3255 mOhm & \textbf{0.8711}\tabularnewline
bga & 2.5253 mOhm & -0.3255 mOhm & \textbf{0.8711}\tabularnewline
bga & 2.5253 mOhm & -0.3255 mOhm & \textbf{0.8711}\tabularnewline
bga & 2.5253 mOhm & -0.3255 mOhm & \textbf{0.8711}\tabularnewline
bga & 2.5253 mOhm & -0.3255 mOhm & \textbf{0.8711}\tabularnewline
bga & 2.5253 mOhm & -0.3255 mOhm & \textbf{0.8711}\tabularnewline
bga & 2.5253 mOhm & -0.3255 mOhm & \textbf{0.8711}\tabularnewline
bga & 2.5253 mOhm & -0.3255 mOhm & \textbf{0.8711}\tabularnewline
bga & 2.5253 mOhm & -0.3255 mOhm & \textbf{0.8711}\tabularnewline
bga & 2.5253 mOhm & -0.3255 mOhm & \textbf{0.8711}\tabularnewline
bga & 2.5253 mOhm & -0.3255 mOhm & \textbf{0.8711}\tabularnewline
bga & 2.5253 mOhm & -0.3255 mOhm & \textbf{0.8711}\tabularnewline
bga & 2.5253 mOhm & -0.3255 mOhm & \textbf{0.8711}\tabularnewline
bga & 2.5253 mOhm & -0.3255 mOhm & \textbf{0.8711}\tabularnewline
bga & 2.5253 mOhm & -0.3255 mOhm & \textbf{0.8711}\tabularnewline
bga & 2.5253 mOhm & -0.3255 mOhm & \textbf{0.8711}\tabularnewline
bga & 2.5253 mOhm & -0.3255 mOhm & \textbf{0.8711}\tabularnewline
bga & 2.5253 mOhm & -0.3255 mOhm & \textbf{0.8711}\tabularnewline
bga & 2.5253 mOhm & -0.3255 mOhm & \textbf{0.8711}\tabularnewline
bga & 2.5253 mOhm & -0.3255 mOhm & \textbf{0.8711}\tabularnewline
bga & 2.5253 mOhm & -0.3255 mOhm & \textbf{0.8711}\tabularnewline
bga & 2.5253 mOhm & -0.3255 mOhm & \textbf{0.8711}\tabularnewline
microbump & 36.8657 mOhm & -2.2667 mOhm & \textbf{0.9385}\tabularnewline
microbump & 36.8657 mOhm & -2.2667 mOhm & \textbf{0.9385}\tabularnewline
microbump & 36.8657 mOhm & -2.2667 mOhm & \textbf{0.9385}\tabularnewline
microbump & 36.8657 mOhm & -2.2667 mOhm & \textbf{0.9385}\tabularnewline
microbump & 36.8657 mOhm & -2.2667 mOhm & \textbf{0.9385}\tabularnewline
microbump & 36.8657 mOhm & -2.2667 mOhm & \textbf{0.9385}\tabularnewline
microbump & 36.8657 mOhm & -2.2667 mOhm & \textbf{0.9385}\tabularnewline
microbump & 36.8657 mOhm & -2.2667 mOhm & \textbf{0.9385}\tabularnewline
microbump & 36.8657 mOhm & -2.2667 mOhm & \textbf{0.9385}\tabularnewline
microbump & 36.8657 mOhm & -2.2667 mOhm & \textbf{0.9385}\tabularnewline
microbump & 36.8657 mOhm & -2.2667 mOhm & \textbf{0.9385}\tabularnewline
microbump & 36.8657 mOhm & -2.2667 mOhm & \textbf{0.9385}\tabularnewline
microbump & 36.8657 mOhm & -2.2667 mOhm & \textbf{0.9385}\tabularnewline
microbump & 36.8657 mOhm & -2.2667 mOhm & \textbf{0.9385}\tabularnewline
microbump & 36.8657 mOhm & -2.2667 mOhm & \textbf{0.9385}\tabularnewline
microbump & 36.8657 mOhm & -2.2667 mOhm & \textbf{0.9385}\tabularnewline
microbump & 36.8657 mOhm & -2.2667 mOhm & \textbf{0.9385}\tabularnewline
microbump & 36.8657 mOhm & -2.2667 mOhm & \textbf{0.9385}\tabularnewline
microbump & 36.8657 mOhm & -2.2667 mOhm & \textbf{0.9385}\tabularnewline
microbump & 36.8657 mOhm & -2.2667 mOhm & \textbf{0.9385}\tabularnewline
microbump & 36.8657 mOhm & -2.2667 mOhm & \textbf{0.9385}\tabularnewline
microbump & 36.8657 mOhm & -2.2667 mOhm & \textbf{0.9385}\tabularnewline
microbump & 36.8657 mOhm & -2.2667 mOhm & \textbf{0.9385}\tabularnewline
microbump & 36.8657 mOhm & -2.2667 mOhm & \textbf{0.9385}\tabularnewline
microbump & 36.8657 mOhm & -2.2667 mOhm & \textbf{0.9385}\tabularnewline
microbump & 36.8657 mOhm & -2.2667 mOhm & \textbf{0.9385}\tabularnewline
microbump & 36.8657 mOhm & -2.2667 mOhm & \textbf{0.9385}\tabularnewline
microbump & 36.8657 mOhm & -2.2667 mOhm & \textbf{0.9385}\tabularnewline
microbump & 36.8657 mOhm & -2.2667 mOhm & \textbf{0.9385}\tabularnewline
microbump & 36.8657 mOhm & -2.2667 mOhm & \textbf{0.9385}\tabularnewline
microbump & 36.8657 mOhm & -2.2667 mOhm & \textbf{0.9385}\tabularnewline
microbump & 36.8657 mOhm & -2.2667 mOhm & \textbf{0.9385}\tabularnewline
microbump & 36.8657 mOhm & -2.2667 mOhm & \textbf{0.9385}\tabularnewline
microbump & 36.8657 mOhm & -2.2667 mOhm & \textbf{0.9385}\tabularnewline
microbump & 36.8657 mOhm & -2.2667 mOhm & \textbf{0.9385}\tabularnewline
microbump & 36.8657 mOhm & -2.2667 mOhm & \textbf{0.9385}\tabularnewline
microbump & 36.8657 mOhm & -2.2667 mOhm & \textbf{0.9385}\tabularnewline
microbump & 36.8657 mOhm & -2.2667 mOhm & \textbf{0.9385}\tabularnewline
microbump & 36.8657 mOhm & -2.2667 mOhm & \textbf{0.9385}\tabularnewline
tsv & 22.2485 mOhm & -2.3440 mOhm & \textbf{0.8946}\tabularnewline
tsv & 22.2485 mOhm & -2.3440 mOhm & \textbf{0.8946}\tabularnewline
tsv & 22.2485 mOhm & -2.3440 mOhm & \textbf{0.8946}\tabularnewline
tsv & 22.2485 mOhm & -2.3440 mOhm & \textbf{0.8946}\tabularnewline
tsv & 22.2485 mOhm & -2.3440 mOhm & \textbf{0.8946}\tabularnewline
tsv & 22.2485 mOhm & -2.3440 mOhm & \textbf{0.8946}\tabularnewline
tsv & 22.2485 mOhm & -2.3440 mOhm & \textbf{0.8946}\tabularnewline
tsv & 22.2485 mOhm & -2.3440 mOhm & \textbf{0.8946}\tabularnewline
tsv & 22.2485 mOhm & -2.3440 mOhm & \textbf{0.8946}\tabularnewline
tsv & 22.2485 mOhm & -2.3440 mOhm & \textbf{0.8946}\tabularnewline
tsv & 22.2485 mOhm & -2.3440 mOhm & \textbf{0.8946}\tabularnewline
tsv & 22.2485 mOhm & -2.3440 mOhm & \textbf{0.8946}\tabularnewline
tsv & 22.2485 mOhm & -2.3440 mOhm & \textbf{0.8946}\tabularnewline
tsv & 22.2485 mOhm & -2.3440 mOhm & \textbf{0.8946}\tabularnewline
tsv & 22.2485 mOhm & -2.3440 mOhm & \textbf{0.8946}\tabularnewline
tsv & 22.2485 mOhm & -2.3440 mOhm & \textbf{0.8946}\tabularnewline
tsv & 22.2485 mOhm & -2.3440 mOhm & \textbf{0.8946}\tabularnewline
tsv & 22.2485 mOhm & -2.3440 mOhm & \textbf{0.8946}\tabularnewline
tsv & 22.2485 mOhm & -2.3440 mOhm & \textbf{0.8946}\tabularnewline
tsv & 22.2485 mOhm & -2.3440 mOhm & \textbf{0.8946}\tabularnewline
tsv & 22.2485 mOhm & -2.3440 mOhm & \textbf{0.8946}\tabularnewline
tsv & 22.2485 mOhm & -2.3440 mOhm & \textbf{0.8946}\tabularnewline
tsv & 22.2485 mOhm & -2.3440 mOhm & \textbf{0.8946}\tabularnewline
tsv & 22.2485 mOhm & -2.3440 mOhm & \textbf{0.8946}\tabularnewline
tsv & 22.2485 mOhm & -2.3440 mOhm & \textbf{0.8946}\tabularnewline
tsv & 22.2485 mOhm & -2.3440 mOhm & \textbf{0.8946}\tabularnewline
tsv & 22.2485 mOhm & -2.3440 mOhm & \textbf{0.8946}\tabularnewline
tsv & 22.2485 mOhm & -2.3440 mOhm & \textbf{0.8946}\tabularnewline
tsv & 22.2485 mOhm & -2.3440 mOhm & \textbf{0.8946}\tabularnewline
tsv & 22.2485 mOhm & -2.3440 mOhm & \textbf{0.8946}\tabularnewline
tsv & 22.2485 mOhm & -2.3440 mOhm & \textbf{0.8946}\tabularnewline
tsv & 22.2485 mOhm & -2.3440 mOhm & \textbf{0.8946}\tabularnewline
tsv & 22.2485 mOhm & -2.3440 mOhm & \textbf{0.8946}\tabularnewline
tsv & 22.2485 mOhm & -2.3440 mOhm & \textbf{0.8946}\tabularnewline
tsv & 22.2485 mOhm & -2.3440 mOhm & \textbf{0.8946}\tabularnewline
tsv & 22.2485 mOhm & -2.3440 mOhm & \textbf{0.8946}\tabularnewline
tsv & 22.2485 mOhm & -2.3440 mOhm & \textbf{0.8946}\tabularnewline
tsv & 22.2485 mOhm & -2.3440 mOhm & \textbf{0.8946}\tabularnewline
tsv & 22.2485 mOhm & -2.3440 mOhm & \textbf{0.8946}\tabularnewline
\bottomrule
\end{longtable}

Against the retired global constant of \textbf{1.326}, applied
identically to all four classes.

\begin{figure}
\centering
\includegraphics{/home/beams0/MGHARIB/Mohamed_stuff/HVR_OPT/results/phase0mr/FIG_0MR_C_transition_decomposition.png}
\caption{Transition decomposition}
\end{figure}

\hypertarget{the-tsv-radius-result-stated-precisely}{%
\subsection{8.2 The TSV radius result, stated
precisely}\label{the-tsv-radius-result-stated-precisely}}

Under the minimum-capability allocation this framework uses,
\texttt{R\_field\ =\ rho\ h\ J\_design\ /\ I} and \textbf{the conductor
cross-section cancels identically} --- not merely the radius. The copper
fraction \texttt{I\ /\ (J\_design\ A\_field)} is likewise independent of
radius, so the effective conductivity, modulus and CTE do not move
either.

The permitted conclusion is therefore \textbf{not} ``TSV radius is
irrelevant''. It is: \emph{the TSV cross-section is not an independent
field-loss lever under minimum-capability provisioning when
\texttt{J\_design} is geometry-independent.} Radius remains decisive for
\textbf{feasibility}:

\begin{longtable}[]{@{}ll@{}}
\toprule
\begin{minipage}[b]{0.47\columnwidth}\raggedright
via height\strut
\end{minipage} & \begin{minipage}[b]{0.47\columnwidth}\raggedright
feasible radii at a 10:1 aspect limit\strut
\end{minipage}\tabularnewline
\midrule
\endhead
\begin{minipage}[t]{0.47\columnwidth}\raggedright
100 um\strut
\end{minipage} & \begin{minipage}[t]{0.47\columnwidth}\raggedright
7.5 to 50 um\strut
\end{minipage}\tabularnewline
\begin{minipage}[t]{0.47\columnwidth}\raggedright
300 um\strut
\end{minipage} & \begin{minipage}[t]{0.47\columnwidth}\raggedright
\textbf{15 to 50 um} --- and 15 um is the reference radius, so the
source geometry sits ON the limit\strut
\end{minipage}\tabularnewline
\begin{minipage}[t]{0.47\columnwidth}\raggedright
1300 um\strut
\end{minipage} & \begin{minipage}[t]{0.47\columnwidth}\raggedright
\textbf{none in the swept range}; the rule
\texttt{r\ \textgreater{}=\ h/20} demands r \textgreater{}= 65 um\strut
\end{minipage}\tabularnewline
\bottomrule
\end{longtable}

\begin{figure}
\centering
\includegraphics{/home/beams0/MGHARIB/Mohamed_stuff/HVR_OPT/results/phase0mr/FIG_0MR_D_tsv_feasibility.png}
\caption{TSV feasibility}
\end{figure}

\newpage

\hypertarget{converter-component-level-partitioning}{%
\section{9. Converter component-level
partitioning}\label{converter-component-level-partitioning}}

The converter is nine components per cell: 3 high-side GaN, 3 low-side
GaN, 2 flying capacitors, 3 inductors and an output capacitor. Their
summed footprint is \textbf{18.01 mm2} against a 500 mm2 die shared by
49 cells --- \textbf{10.20 mm2 per cell}. A single-level VR therefore
misses the die shadow by \textbf{1.76x}, and vertical partitioning is a
precondition rather than an optimization.

Each template is checked on six independent questions, not on area
alone:

\begin{longtable}[]{@{}llllllll@{}}
\toprule
\begin{minipage}[b]{0.10\columnwidth}\raggedright
template\strut
\end{minipage} & \begin{minipage}[b]{0.10\columnwidth}\raggedright
tiers\strut
\end{minipage} & \begin{minipage}[b]{0.10\columnwidth}\raggedright
charged height / available\strut
\end{minipage} & \begin{minipage}[b]{0.10\columnwidth}\raggedright
area\strut
\end{minipage} & \begin{minipage}[b]{0.10\columnwidth}\raggedright
height\strut
\end{minipage} & \begin{minipage}[b]{0.10\columnwidth}\raggedright
routing\strut
\end{minipage} & \begin{minipage}[b]{0.10\columnwidth}\raggedright
thermal\strut
\end{minipage} & \begin{minipage}[b]{0.10\columnwidth}\raggedright
clearance\strut
\end{minipage}\tabularnewline
\midrule
\endhead
\begin{minipage}[t]{0.10\columnwidth}\raggedright
T1\_EMBEDDED\_4LAYER in Z1\_EMBEDDED\strut
\end{minipage} & \begin{minipage}[t]{0.10\columnwidth}\raggedright
4\strut
\end{minipage} & \begin{minipage}[t]{0.10\columnwidth}\raggedright
1226 / 1250 um\strut
\end{minipage} & \begin{minipage}[t]{0.10\columnwidth}\raggedright
True\strut
\end{minipage} & \begin{minipage}[t]{0.10\columnwidth}\raggedright
True\strut
\end{minipage} & \begin{minipage}[t]{0.10\columnwidth}\raggedright
False\strut
\end{minipage} & \begin{minipage}[t]{0.10\columnwidth}\raggedright
False\strut
\end{minipage} & \begin{minipage}[t]{0.10\columnwidth}\raggedright
OPEN\strut
\end{minipage}\tabularnewline
\begin{minipage}[t]{0.10\columnwidth}\raggedright
T1\_EMBEDDED\_4LAYER in Z2\_BACKSIDE\strut
\end{minipage} & \begin{minipage}[t]{0.10\columnwidth}\raggedright
4\strut
\end{minipage} & \begin{minipage}[t]{0.10\columnwidth}\raggedright
1921 / 1100 um\strut
\end{minipage} & \begin{minipage}[t]{0.10\columnwidth}\raggedright
True\strut
\end{minipage} & \begin{minipage}[t]{0.10\columnwidth}\raggedright
False\strut
\end{minipage} & \begin{minipage}[t]{0.10\columnwidth}\raggedright
False\strut
\end{minipage} & \begin{minipage}[t]{0.10\columnwidth}\raggedright
True\strut
\end{minipage} & \begin{minipage}[t]{0.10\columnwidth}\raggedright
OPEN\strut
\end{minipage}\tabularnewline
\begin{minipage}[t]{0.10\columnwidth}\raggedright
T1\_EMBEDDED\_4LAYER in Z4\_POWER\_DIE\strut
\end{minipage} & \begin{minipage}[t]{0.10\columnwidth}\raggedright
4\strut
\end{minipage} & \begin{minipage}[t]{0.10\columnwidth}\raggedright
1921 / 1200 um\strut
\end{minipage} & \begin{minipage}[t]{0.10\columnwidth}\raggedright
True\strut
\end{minipage} & \begin{minipage}[t]{0.10\columnwidth}\raggedright
False\strut
\end{minipage} & \begin{minipage}[t]{0.10\columnwidth}\raggedright
False\strut
\end{minipage} & \begin{minipage}[t]{0.10\columnwidth}\raggedright
False\strut
\end{minipage} & \begin{minipage}[t]{0.10\columnwidth}\raggedright
OPEN\strut
\end{minipage}\tabularnewline
\begin{minipage}[t]{0.10\columnwidth}\raggedright
T4\_SPLIT\_SC\_BUCK in Z1\_EMBEDDED\strut
\end{minipage} & \begin{minipage}[t]{0.10\columnwidth}\raggedright
2\strut
\end{minipage} & \begin{minipage}[t]{0.10\columnwidth}\raggedright
1301 / 1250 um\strut
\end{minipage} & \begin{minipage}[t]{0.10\columnwidth}\raggedright
False\strut
\end{minipage} & \begin{minipage}[t]{0.10\columnwidth}\raggedright
False\strut
\end{minipage} & \begin{minipage}[t]{0.10\columnwidth}\raggedright
True\strut
\end{minipage} & \begin{minipage}[t]{0.10\columnwidth}\raggedright
True\strut
\end{minipage} & \begin{minipage}[t]{0.10\columnwidth}\raggedright
OPEN\strut
\end{minipage}\tabularnewline
\begin{minipage}[t]{0.10\columnwidth}\raggedright
T4\_SPLIT\_SC\_BUCK in Z2\_BACKSIDE\strut
\end{minipage} & \begin{minipage}[t]{0.10\columnwidth}\raggedright
2\strut
\end{minipage} & \begin{minipage}[t]{0.10\columnwidth}\raggedright
1301 / 1100 um\strut
\end{minipage} & \begin{minipage}[t]{0.10\columnwidth}\raggedright
False\strut
\end{minipage} & \begin{minipage}[t]{0.10\columnwidth}\raggedright
False\strut
\end{minipage} & \begin{minipage}[t]{0.10\columnwidth}\raggedright
True\strut
\end{minipage} & \begin{minipage}[t]{0.10\columnwidth}\raggedright
True\strut
\end{minipage} & \begin{minipage}[t]{0.10\columnwidth}\raggedright
OPEN\strut
\end{minipage}\tabularnewline
\begin{minipage}[t]{0.10\columnwidth}\raggedright
T4\_SPLIT\_SC\_BUCK in Z4\_POWER\_DIE\strut
\end{minipage} & \begin{minipage}[t]{0.10\columnwidth}\raggedright
2\strut
\end{minipage} & \begin{minipage}[t]{0.10\columnwidth}\raggedright
1301 / 1200 um\strut
\end{minipage} & \begin{minipage}[t]{0.10\columnwidth}\raggedright
False\strut
\end{minipage} & \begin{minipage}[t]{0.10\columnwidth}\raggedright
False\strut
\end{minipage} & \begin{minipage}[t]{0.10\columnwidth}\raggedright
True\strut
\end{minipage} & \begin{minipage}[t]{0.10\columnwidth}\raggedright
True\strut
\end{minipage} & \begin{minipage}[t]{0.10\columnwidth}\raggedright
OPEN\strut
\end{minipage}\tabularnewline
\begin{minipage}[t]{0.10\columnwidth}\raggedright
T2\_BACKSIDE\_CHIPLET in Z1\_EMBEDDED\strut
\end{minipage} & \begin{minipage}[t]{0.10\columnwidth}\raggedright
3\strut
\end{minipage} & \begin{minipage}[t]{0.10\columnwidth}\raggedright
1305 / 1250 um\strut
\end{minipage} & \begin{minipage}[t]{0.10\columnwidth}\raggedright
False\strut
\end{minipage} & \begin{minipage}[t]{0.10\columnwidth}\raggedright
False\strut
\end{minipage} & \begin{minipage}[t]{0.10\columnwidth}\raggedright
False\strut
\end{minipage} & \begin{minipage}[t]{0.10\columnwidth}\raggedright
True\strut
\end{minipage} & \begin{minipage}[t]{0.10\columnwidth}\raggedright
OPEN\strut
\end{minipage}\tabularnewline
\begin{minipage}[t]{0.10\columnwidth}\raggedright
T2\_BACKSIDE\_CHIPLET in Z2\_BACKSIDE\strut
\end{minipage} & \begin{minipage}[t]{0.10\columnwidth}\raggedright
3\strut
\end{minipage} & \begin{minipage}[t]{0.10\columnwidth}\raggedright
1305 / 1100 um\strut
\end{minipage} & \begin{minipage}[t]{0.10\columnwidth}\raggedright
False\strut
\end{minipage} & \begin{minipage}[t]{0.10\columnwidth}\raggedright
False\strut
\end{minipage} & \begin{minipage}[t]{0.10\columnwidth}\raggedright
False\strut
\end{minipage} & \begin{minipage}[t]{0.10\columnwidth}\raggedright
True\strut
\end{minipage} & \begin{minipage}[t]{0.10\columnwidth}\raggedright
OPEN\strut
\end{minipage}\tabularnewline
\begin{minipage}[t]{0.10\columnwidth}\raggedright
T2\_BACKSIDE\_CHIPLET in Z4\_POWER\_DIE\strut
\end{minipage} & \begin{minipage}[t]{0.10\columnwidth}\raggedright
3\strut
\end{minipage} & \begin{minipage}[t]{0.10\columnwidth}\raggedright
1305 / 1200 um\strut
\end{minipage} & \begin{minipage}[t]{0.10\columnwidth}\raggedright
False\strut
\end{minipage} & \begin{minipage}[t]{0.10\columnwidth}\raggedright
False\strut
\end{minipage} & \begin{minipage}[t]{0.10\columnwidth}\raggedright
False\strut
\end{minipage} & \begin{minipage}[t]{0.10\columnwidth}\raggedright
True\strut
\end{minipage} & \begin{minipage}[t]{0.10\columnwidth}\raggedright
OPEN\strut
\end{minipage}\tabularnewline
\begin{minipage}[t]{0.10\columnwidth}\raggedright
T0\_SINGLE\_LEVEL in Z1\_EMBEDDED\strut
\end{minipage} & \begin{minipage}[t]{0.10\columnwidth}\raggedright
1\strut
\end{minipage} & \begin{minipage}[t]{0.10\columnwidth}\raggedright
685 / 1250 um\strut
\end{minipage} & \begin{minipage}[t]{0.10\columnwidth}\raggedright
False\strut
\end{minipage} & \begin{minipage}[t]{0.10\columnwidth}\raggedright
True\strut
\end{minipage} & \begin{minipage}[t]{0.10\columnwidth}\raggedright
True\strut
\end{minipage} & \begin{minipage}[t]{0.10\columnwidth}\raggedright
True\strut
\end{minipage} & \begin{minipage}[t]{0.10\columnwidth}\raggedright
OPEN\strut
\end{minipage}\tabularnewline
\begin{minipage}[t]{0.10\columnwidth}\raggedright
T0\_SINGLE\_LEVEL in Z2\_BACKSIDE\strut
\end{minipage} & \begin{minipage}[t]{0.10\columnwidth}\raggedright
1\strut
\end{minipage} & \begin{minipage}[t]{0.10\columnwidth}\raggedright
685 / 1100 um\strut
\end{minipage} & \begin{minipage}[t]{0.10\columnwidth}\raggedright
False\strut
\end{minipage} & \begin{minipage}[t]{0.10\columnwidth}\raggedright
True\strut
\end{minipage} & \begin{minipage}[t]{0.10\columnwidth}\raggedright
True\strut
\end{minipage} & \begin{minipage}[t]{0.10\columnwidth}\raggedright
True\strut
\end{minipage} & \begin{minipage}[t]{0.10\columnwidth}\raggedright
OPEN\strut
\end{minipage}\tabularnewline
\begin{minipage}[t]{0.10\columnwidth}\raggedright
T0\_SINGLE\_LEVEL in Z4\_POWER\_DIE\strut
\end{minipage} & \begin{minipage}[t]{0.10\columnwidth}\raggedright
1\strut
\end{minipage} & \begin{minipage}[t]{0.10\columnwidth}\raggedright
685 / 1200 um\strut
\end{minipage} & \begin{minipage}[t]{0.10\columnwidth}\raggedright
False\strut
\end{minipage} & \begin{minipage}[t]{0.10\columnwidth}\raggedright
True\strut
\end{minipage} & \begin{minipage}[t]{0.10\columnwidth}\raggedright
True\strut
\end{minipage} & \begin{minipage}[t]{0.10\columnwidth}\raggedright
True\strut
\end{minipage} & \begin{minipage}[t]{0.10\columnwidth}\raggedright
OPEN\strut
\end{minipage}\tabularnewline
\bottomrule
\end{longtable}

\textbf{No template is fully realizable.} The closest is the embedded
four-layer stack in the embedded architecture, whose charged height of
1226 um reproduces an independently derived earlier figure exactly and
fits the 1250 um available. It is blocked on routing and on thermal
depth.

\textbf{Three package clearances have no source} --- die-to-component
gap, cavity sidewall allowance, cooled-surface standoff --- so every
template carries CLEARANCE OPEN and \textbf{none may be called
buildable}. Treating them as zero would convert three unknowns into a
pass.

\hypertarget{the-intermediate-rail}{%
\subsection{9.1 The intermediate rail}\label{the-intermediate-rail}}

A partition that splits by VOLTAGE DOMAIN puts the 16 V intermediate
rail across the tier boundary, carrying \textbf{1.44 A where the 1 V
rail would carry 20.4 A} --- a factor of 200 in \texttt{I\^{}2}. That is
why the split SC/buck template has essentially zero vertical intra-VR
loss.

\begin{figure}
\centering
\includegraphics{/home/beams0/MGHARIB/Mohamed_stuff/HVR_OPT/results/phase0mr/FIG_0MR_B_component_power_path.png}
\caption{Component power path}
\end{figure}

\newpage

\hypertarget{intra-vr-routing}{%
\section{10. Intra-VR routing}\label{intra-vr-routing}}

\textbf{The metal between the converter's own components is now in the
loss total.} Before Phase 0M-R it was absent entirely: the framework
modelled VR-to-package and package-to-POL routing and nothing inside the
converter.

There is no double counting, and the reason is structural rather than a
comment. The upstream converter model owns every device and passive
parasitic --- switch \texttt{R\_DS(on)}, inductor DCR, capacitor ESR,
and all switching loss --- and contains no metallization. The ledger
enforces \textbf{one owner per watt}: a second writer to any term
raises, and the conduction split is checked against the upstream
aggregate to machine precision before a total may be produced.

\begin{longtable}[]{@{}ll@{}}
\toprule
\begin{minipage}[b]{0.47\columnwidth}\raggedright
term\strut
\end{minipage} & \begin{minipage}[b]{0.47\columnwidth}\raggedright
owner\strut
\end{minipage}\tabularnewline
\midrule
\endhead
\begin{minipage}[t]{0.47\columnwidth}\raggedright
\texttt{P\_switch\_hs\_conduction},
\texttt{P\_switch\_ls\_conduction}\strut
\end{minipage} & \begin{minipage}[t]{0.47\columnwidth}\raggedright
upstream conduction formula\strut
\end{minipage}\tabularnewline
\begin{minipage}[t]{0.47\columnwidth}\raggedright
\texttt{P\_inductor\_DCR}, \texttt{P\_cap\_ESR}\strut
\end{minipage} & \begin{minipage}[t]{0.47\columnwidth}\raggedright
upstream conduction formula\strut
\end{minipage}\tabularnewline
\begin{minipage}[t]{0.47\columnwidth}\raggedright
\texttt{P\_switching\_hard}, \texttt{P\_gate}, \texttt{P\_Coss}\strut
\end{minipage} & \begin{minipage}[t]{0.47\columnwidth}\raggedright
upstream switching formula\strut
\end{minipage}\tabularnewline
\begin{minipage}[t]{0.47\columnwidth}\raggedright
\texttt{P\_intraVR\_metal}, \texttt{P\_intraVR\_vertical}\strut
\end{minipage} & \begin{minipage}[t]{0.47\columnwidth}\raggedright
\textbf{this framework, new}\strut
\end{minipage}\tabularnewline
\begin{minipage}[t]{0.47\columnwidth}\raggedright
\texttt{P\_TSV}, \texttt{P\_bumps}\strut
\end{minipage} & \begin{minipage}[t]{0.47\columnwidth}\raggedright
vertical tier model, \texttt{R\_bulk\ +\ dR}\strut
\end{minipage}\tabularnewline
\begin{minipage}[t]{0.47\columnwidth}\raggedright
\texttt{P\_pkg\_RDL}\strut
\end{minipage} & \begin{minipage}[t]{0.47\columnwidth}\raggedright
multilayer resistive sheet\strut
\end{minipage}\tabularnewline
\bottomrule
\end{longtable}

The routing model is \textbf{template-aware}: components on the same
tier need a lateral run, components on different tiers need a via stack,
and the via-stack height is set by the components below it. Phase 0M's
claim that intra-VR routing is common-mode across architectures held
only while every architecture used the same layout; measured across
partition templates it varies by \textbf{7.5 \%}.

\newpage

\hypertarget{the-cooling-model}{%
\section{11. The cooling model}\label{the-cooling-model}}

\textbf{The cooled boundary is derived from hardware, not calibrated to
a target.}

\[-k \nabla T \cdot n = h(x)\left[T_s(x) - T_f(x)\right], \qquad T_f(x) = T_{in} + \frac{Q(0 \to x)}{\dot{m} c_p}\]

\texttt{h} comes from a Nusselt correlation for the channel geometry,
corrected by fin efficiency and area amplification. \texttt{T\_f} comes
from the heat actually removed upstream of each face, integrated from
the solved flux and iterated with the wall temperature to a fixed point.
The design variables are coolant, inlet temperature, mass flow, channel
width, height and pitch, base thickness and solid conductivity.

\begin{longtable}[]{@{}lllllllll@{}}
\toprule
\begin{minipage}[b]{0.08\columnwidth}\raggedright
design\strut
\end{minipage} & \begin{minipage}[b]{0.08\columnwidth}\raggedright
m\_dot\strut
\end{minipage} & \begin{minipage}[b]{0.08\columnwidth}\raggedright
Re\strut
\end{minipage} & \begin{minipage}[b]{0.08\columnwidth}\raggedright
regime\strut
\end{minipage} & \begin{minipage}[b]{0.08\columnwidth}\raggedright
h\_effective\strut
\end{minipage} & \begin{minipage}[b]{0.08\columnwidth}\raggedright
coolant rise\strut
\end{minipage} & \begin{minipage}[b]{0.08\columnwidth}\raggedright
T\_j\strut
\end{minipage} & \begin{minipage}[b]{0.08\columnwidth}\raggedright
dp\strut
\end{minipage} & \begin{minipage}[b]{0.08\columnwidth}\raggedright
pumping\strut
\end{minipage}\tabularnewline
\midrule
\endhead
\begin{minipage}[t]{0.08\columnwidth}\raggedright
TOP\_SI\_MICROCHANNEL\_27\strut
\end{minipage} & \begin{minipage}[t]{0.08\columnwidth}\raggedright
27.0 g/s\strut
\end{minipage} & \begin{minipage}[t]{0.08\columnwidth}\raggedright
1265\strut
\end{minipage} & \begin{minipage}[t]{0.08\columnwidth}\raggedright
laminar\strut
\end{minipage} & \begin{minipage}[t]{0.08\columnwidth}\raggedright
115553\strut
\end{minipage} & \begin{minipage}[t]{0.08\columnwidth}\raggedright
9.97 K\strut
\end{minipage} & \begin{minipage}[t]{0.08\columnwidth}\raggedright
44.5 C\strut
\end{minipage} & \begin{minipage}[t]{0.08\columnwidth}\raggedright
538 kPa\strut
\end{minipage} & \begin{minipage}[t]{0.08\columnwidth}\raggedright
14.70 W\strut
\end{minipage}\tabularnewline
\begin{minipage}[t]{0.08\columnwidth}\raggedright
TOP\_SI\_MICROCHANNEL\_50\strut
\end{minipage} & \begin{minipage}[t]{0.08\columnwidth}\raggedright
50.0 g/s\strut
\end{minipage} & \begin{minipage}[t]{0.08\columnwidth}\raggedright
2342\strut
\end{minipage} & \begin{minipage}[t]{0.08\columnwidth}\raggedright
transitional\strut
\end{minipage} & \begin{minipage}[t]{0.08\columnwidth}\raggedright
215115\strut
\end{minipage} & \begin{minipage}[t]{0.08\columnwidth}\raggedright
5.38 K\strut
\end{minipage} & \begin{minipage}[t]{0.08\columnwidth}\raggedright
35.5 C\strut
\end{minipage} & \begin{minipage}[t]{0.08\columnwidth}\raggedright
1463 kPa\strut
\end{minipage} & \begin{minipage}[t]{0.08\columnwidth}\raggedright
74.06 W\strut
\end{minipage}\tabularnewline
\begin{minipage}[t]{0.08\columnwidth}\raggedright
TOP\_CU\_COLDPLATE\_27\strut
\end{minipage} & \begin{minipage}[t]{0.08\columnwidth}\raggedright
27.0 g/s\strut
\end{minipage} & \begin{minipage}[t]{0.08\columnwidth}\raggedright
1482\strut
\end{minipage} & \begin{minipage}[t]{0.08\columnwidth}\raggedright
laminar\strut
\end{minipage} & \begin{minipage}[t]{0.08\columnwidth}\raggedright
70478\strut
\end{minipage} & \begin{minipage}[t]{0.08\columnwidth}\raggedright
9.97 K\strut
\end{minipage} & \begin{minipage}[t]{0.08\columnwidth}\raggedright
56.9 C\strut
\end{minipage} & \begin{minipage}[t]{0.08\columnwidth}\raggedright
83 kPa\strut
\end{minipage} & \begin{minipage}[t]{0.08\columnwidth}\raggedright
2.26 W\strut
\end{minipage}\tabularnewline
\begin{minipage}[t]{0.08\columnwidth}\raggedright
TWO\_SIDED\_SI\_14\strut
\end{minipage} & \begin{minipage}[t]{0.08\columnwidth}\raggedright
13.5 g/s\strut
\end{minipage} & \begin{minipage}[t]{0.08\columnwidth}\raggedright
632\strut
\end{minipage} & \begin{minipage}[t]{0.08\columnwidth}\raggedright
laminar\strut
\end{minipage} & \begin{minipage}[t]{0.08\columnwidth}\raggedright
92009\strut
\end{minipage} & \begin{minipage}[t]{0.08\columnwidth}\raggedright
9.97 K\strut
\end{minipage} & \begin{minipage}[t]{0.08\columnwidth}\raggedright
37.2 C\strut
\end{minipage} & \begin{minipage}[t]{0.08\columnwidth}\raggedright
269 kPa\strut
\end{minipage} & \begin{minipage}[t]{0.08\columnwidth}\raggedright
7.35 W\strut
\end{minipage}\tabularnewline
\begin{minipage}[t]{0.08\columnwidth}\raggedright
SUBSTRATE\_EMBEDDED\_27\strut
\end{minipage} & \begin{minipage}[t]{0.08\columnwidth}\raggedright
27.0 g/s\strut
\end{minipage} & \begin{minipage}[t]{0.08\columnwidth}\raggedright
2965\strut
\end{minipage} & \begin{minipage}[t]{0.08\columnwidth}\raggedright
transitional\strut
\end{minipage} & \begin{minipage}[t]{0.08\columnwidth}\raggedright
125399\strut
\end{minipage} & \begin{minipage}[t]{0.08\columnwidth}\raggedright
9.97 K\strut
\end{minipage} & \begin{minipage}[t]{0.08\columnwidth}\raggedright
42.9 C\strut
\end{minipage} & \begin{minipage}[t]{0.08\columnwidth}\raggedright
574 kPa\strut
\end{minipage} & \begin{minipage}[t]{0.08\columnwidth}\raggedright
15.69 W\strut
\end{minipage}\tabularnewline
\bottomrule
\end{longtable}

\textbf{Pumping power is reported and never folded into the electrical
total.} \texttt{P\_electronics} and \texttt{P\_pumping} stay separate;
pump efficiency is unsourced, so hydraulic power is primary and a
35/55/75 \% shaft-power band is carried beside it.

\hypertarget{what-this-replaced-and-why-it-mattered}{%
\subsection{11.1 What this replaced, and why it
mattered}\label{what-this-replaced-and-why-it-mattered}}

The previous boundary was a single inverse-calibrated coefficient,
\textbf{h = 166 000 W/m2K}, with one uniform sink temperature, chosen so
the package would reach a target junction temperature. Three things were
hidden inside it:

\begin{itemize}
\tightlist
\item
  \textbf{no laminar microchannel architecture reaches that
  coefficient.} The best gets 0.70x. Reaching it needs 50 g/s in the
  transitional regime, \textbf{74 W of pumping --- 7.4 \% of the load
  ---} and 1.46 MPa;
\item
  at the flow rate the reference system states, the coolant would rise
  \textbf{168 K}, which single-phase water cannot do. A 10 K rise needs
  26.9 g/s;
\item
  a single sink temperature cannot produce a streamwise gradient, and a
  real one moves the die warpage by several per cent.
\end{itemize}

\begin{figure}
\centering
\includegraphics{/home/beams0/MGHARIB/Mohamed_stuff/HVR_OPT/results/phase0mr/FIG_0MR_F_cooling.png}
\caption{Cooling}
\end{figure}

\hypertarget{validation}{%
\subsection{11.2 Validation}\label{validation}}

The compact model has three parts and each is validated by a tool that
can pose it:

\begin{longtable}[]{@{}lll@{}}
\toprule
part & method & result\tabularnewline
\midrule
\endhead
fluid side & analytic energy balance & exact for every
design\tabularnewline
solid side & MOOSE on the explicit finned geometry & see
below\tabularnewline
coupling & fixed-point iteration, energy checked & residual reported per
solve\tabularnewline
\bottomrule
\end{longtable}

\begin{longtable}[]{@{}llllll@{}}
\toprule
case & fin efficiency & R compact & R MOOSE & error &
verdict\tabularnewline
\midrule
\endhead
TuckermanPease\_50x302 & 0.7666 & 6.4033e-06 & 6.1699e-06 & 3.78 \% &
\textbf{ACCEPTED}\tabularnewline
wide\_100x200 & 0.8562 & 6.1063e-06 & 5.5458e-06 & 10.11 \% &
\textbf{REJECTED}\tabularnewline
tall\_50x500 & 0.5632 & 5.3826e-06 & 5.4517e-06 & 1.27 \% &
\textbf{ACCEPTED}\tabularnewline
copper\_100x500 & 0.9322 & 1.1633e-05 & 1.0971e-05 & 6.04 \% &
\textbf{CONDITIONAL}\tabularnewline
\bottomrule
\end{longtable}

\textbf{The compact solid model is valid for tall thin fins and degrades
for short thick ones}, which is the expected failure of a 1-D fin
assumption. The production geometry is inside the accepted band at 3.78
\%; the short-fin case is REJECTED at 10.11 \% against a band declared
before the runs. That boundary is recorded, not glossed.

\newpage

\hypertarget{thermal-homogenization}{%
\section{12. Thermal homogenization}\label{thermal-homogenization}}

Every heterogeneous layer carries an anisotropic effective tensor,
validated against explicit-geometry finite volumes on unit cells with a
predeclared 10 \% / 25 \% acceptance band. The unit cell is its own
periodic cell: for a centred inclusion in a square cell, symmetry puts
zero normal flux on every mid-plane, and a 1x1 against 2x2 check
confirms it to 0.000 \%.

\begin{longtable}[]{@{}lll@{}}
\toprule
\begin{minipage}[b]{0.30\columnwidth}\raggedright
layer\strut
\end{minipage} & \begin{minipage}[b]{0.30\columnwidth}\raggedright
what survives\strut
\end{minipage} & \begin{minipage}[b]{0.30\columnwidth}\raggedright
what does not\strut
\end{minipage}\tabularnewline
\midrule
\endhead
\begin{minipage}[t]{0.30\columnwidth}\raggedright
TSV + Si + liner\strut
\end{minipage} & \begin{minipage}[t]{0.30\columnwidth}\raggedright
Voigt for k\_z, better than 1 \%\strut
\end{minipage} & \begin{minipage}[t]{0.30\columnwidth}\raggedright
the pre-existing framework's transverse model, 19 to 74 \% low, because
it puts the liner in series with the copper\strut
\end{minipage}\tabularnewline
\begin{minipage}[t]{0.30\columnwidth}\raggedright
bump + underfill\strut
\end{minipage} & \begin{minipage}[t]{0.30\columnwidth}\raggedright
Voigt for k\_z and Maxwell for k\_xy, both under 1.5 \%\strut
\end{minipage} & \begin{minipage}[t]{0.30\columnwidth}\raggedright
the TSV formula, which must not be reused: the contrast and the matrix
are different\strut
\end{minipage}\tabularnewline
\begin{minipage}[t]{0.30\columnwidth}\raggedright
RDL stack\strut
\end{minipage} & \begin{minipage}[t]{0.30\columnwidth}\raggedright
direction-resolved tensor per stack organisation\strut
\end{minipage} & \begin{minipage}[t]{0.30\columnwidth}\raggedright
any single isotropic value\strut
\end{minipage}\tabularnewline
\bottomrule
\end{longtable}

\hypertarget{the-rdl-where-the-number-depends-on-the-organisation}{%
\subsection{12.1 The RDL, where the number depends on the
organisation}\label{the-rdl-where-the-number-depends-on-the-organisation}}

\begin{longtable}[]{@{}llllll@{}}
\toprule
organisation & k\_x & k\_y & k\_z & in-plane x/y & in-plane /
through\tabularnewline
\midrule
\endhead
LINES\_1D & 233.47 & 1.76 & 1.24 & \textbf{132.55} & 189x\tabularnewline
ORTHOGONAL & 117.63 & 117.63 & 1.24 & \textbf{1.00} & 95x\tabularnewline
MESH & 260.74 & 260.74 & 1.27 & \textbf{1.00} & 205x\tabularnewline
PLANE & 280.10 & 280.10 & 1.28 & \textbf{1.00} & 219x\tabularnewline
\bottomrule
\end{longtable}

A single-direction line layer is strongly anisotropic in plane.
\textbf{Every realistic organisation --- orthogonal routing, a mesh, a
plane --- is in-plane isotropic to 1.00.} The real anisotropy is
through-thickness against in-plane, 57x to 270x, because the inter-layer
dielectric sits in series through the stack while the copper is in
parallel across it.

That matters because the production package still gives this layer
\textbf{one isotropic 0.7 W/m.K}: roughly right through thickness by
accident, and two to three orders of magnitude too low in plane. The
framework is missing a major lateral heat-spreading path, and the
correction is built and measured but is not yet the default.

\begin{figure}
\centering
\includegraphics{/home/beams0/MGHARIB/Mohamed_stuff/HVR_OPT/results/phase0mr/FIG_0MR_E_rdl_tensor.png}
\caption{RDL tensor}
\end{figure}

\hypertarget{what-is-deliberately-not-homogenized}{%
\subsection{12.2 What is deliberately NOT
homogenized}\label{what-is-deliberately-not-homogenized}}

The local solder and transition regions keep explicit geometry.
Homogenizing them would destroy the gradients the reliability chain
reads. Global/local submodelling is retained throughout: tensors in the
package model, explicit geometry in the joint and transition submodels.

\newpage

\hypertarget{the-mechanical-model}{%
\section{13. The mechanical model}\label{the-mechanical-model}}

\hypertarget{global}{%
\subsection{13.1 Global}\label{global}}

A linear thermoelastic finite-element model of the whole package: board,
BGA field, converter layer, substrate with a homogenized TSV region,
interface layer, die and TIM. Silicon is cubic-anisotropic; everything
else is isotropic. The converged temperature field from Section 11
drives it.

It produces two things the rest of the chain uses: \textbf{die and
package warpage}, and \textbf{the relative in-plane displacement across
the corner BGA standoff} --- the quantity handed to the local solder
model.

\hypertarget{local-solder-submodel}{%
\subsection{13.2 Local solder submodel}\label{local-solder-submodel}}

A graded finite-element model of the single worst joint, driven by the
corner displacement, cycled to a stabilised response. The constitutive
law is the overstress form in Section 4.3. The metric is the accumulated
equivalent inelastic strain per stabilised cycle, volume-averaged over
the solder core with a one-element corner exclusion band and matched to
the averaging volume of the life law it feeds.

\hypertarget{the-rdl-in-the-mechanical-model}{%
\subsection{13.3 The RDL in the mechanical
model}\label{the-rdl-in-the-mechanical-model}}

\textbf{The redistribution metal is not in the mechanical mesh.} The
layer between the substrate and the die is 10 um of homogenized
microbump and underfill at 8 GPa; the real RDL stack is 48 to 104 GPa in
plane. A fidelity ladder was built and measured:

\begin{longtable}[]{@{}lllll@{}}
\toprule
rung & interface layer & Z1 corner shear & Z2 corner shear &
Z1/Z2\tabularnewline
\midrule
\endhead
M-RDL0 & as built & 2.7486 um & 2.0796 um &
\textbf{1.3217}\tabularnewline
M-RDL1 & composite, 10 um & 2.6502 um & 2.0044 um &
\textbf{1.3222}\tabularnewline
M-RDL2 & composite, 50 um & 2.9358 um & 2.2237 um &
\textbf{1.3202}\tabularnewline
\bottomrule
\end{longtable}

Adding the RDL moves the absolute corner shear by \textbf{-3.6 \% to
+6.8 \%} and the architecture RATIO by at most \textbf{0.11 \%}. It
changes the numbers and it does not change the ordering. \textbf{No RDL
allowable is invented}, so the new stress remains response coupling and
not a constraint.

\newpage

\hypertarget{the-solder-reliability-model-and-what-it-cannot-do}{%
\section{14. The solder reliability model, and what it cannot
do}\label{the-solder-reliability-model-and-what-it-cannot-do}}

\hypertarget{what-is-solid}{%
\subsection{14.1 What is solid}\label{what-is-solid}}

\begin{itemize}
\tightlist
\item
  The metric is convergent, mesh-independent within its declared band,
  and verified to equal the solver's own accumulated inelastic strain
  including creep.
\item
  Four SnAgCu fatigue formulations are admitted for RELATIVE life, their
  primary sources obtained and one calibration reproduced from the
  authors' own corrected coefficients.
\item
  The inelastic energy is reconstructed from the solver's mechanical
  strain, so no separately parsed temperature enters the energy path.
\end{itemize}

\hypertarget{what-is-conditional-the-dominant-open-item}{%
\subsection{14.2 What is conditional --- the dominant open
item}\label{what-is-conditional-the-dominant-open-item}}

\begin{quote}
The creep half of the solder card is fitted on \textbf{Sn3.9Ag0.6Cu, not
SAC305}, over \textbf{7.51 to 15.92 MPa}, while the joint's
inelastic-strain-weighted median stress is about \textbf{82 MPa} and
\textbf{99.8 \% of the accumulated inelastic strain is generated outside
the fitted window}. The weight is the solver's own dPEEQ times element
volume; no damage law is applied, so this is a RESPONSE weighting.
\end{quote}

Two consequences that must not be conflated:

\begin{itemize}
\tightlist
\item
  \textbf{The ORDERING is robust.} It originates in the global
  thermoelastic model, which imposes a 32 \% larger corner shear on Z1
  and contains no solder law at all. The global package supplies a
  constitutive-law-independent deformation driver favoring Z2, and every
  presently testable constitutive and life-law combination preserves
  that ordering.
\item
  \textbf{The MAGNITUDE is a conditional envelope of 1.19 to 1.22}, not
  a validated number. The permitted wording is ``conditional
  relative-life envelope''. ``Validated relative life'' is not earned.
\end{itemize}

\hypertarget{what-is-absent}{%
\subsection{14.3 What is absent}\label{what-is-absent}}

\textbf{Absolute life.} No required life exists anywhere in the sources,
and the absolute life constant cannot be calibrated on top of an
extrapolated constitutive model without burying the constitutive error
inside the constant.

\textbf{A mechanical allowable.} Nine candidate thresholds were tested
against six admissibility conditions and none passed. Without one, the
reliability response cannot become a feasibility constraint, and the
falsification gate that would test that remains untestable.

\hypertarget{which-mechanisms-exist-at-all}{%
\subsection{14.4 Which mechanisms exist at
all}\label{which-mechanisms-exist-at-all}}

\begin{longtable}[]{@{}ll@{}}
\toprule
\begin{minipage}[b]{0.47\columnwidth}\raggedright
mechanism\strut
\end{minipage} & \begin{minipage}[b]{0.47\columnwidth}\raggedright
status\strut
\end{minipage}\tabularnewline
\midrule
\endhead
\begin{minipage}[t]{0.47\columnwidth}\raggedright
global bending\strut
\end{minipage} & \begin{minipage}[t]{0.47\columnwidth}\raggedright
modelled --- a driver, not a failure\strut
\end{minipage}\tabularnewline
\begin{minipage}[t]{0.47\columnwidth}\raggedright
solder fatigue\strut
\end{minipage} & \begin{minipage}[t]{0.47\columnwidth}\raggedright
\textbf{modelled}, convergent, conditional magnitude\strut
\end{minipage}\tabularnewline
\begin{minipage}[t]{0.47\columnwidth}\raggedright
solder creep / unified viscoplasticity\strut
\end{minipage} & \begin{minipage}[t]{0.47\columnwidth}\raggedright
modelled\strut
\end{minipage}\tabularnewline
\begin{minipage}[t]{0.47\columnwidth}\raggedright
rate-independent plasticity\strut
\end{minipage} & \begin{minipage}[t]{0.47\columnwidth}\raggedright
not separately modelled, and not justified on top of a unified law\strut
\end{minipage}\tabularnewline
\begin{minipage}[t]{0.47\columnwidth}\raggedright
organic board and underfill viscoelasticity\strut
\end{minipage} & \begin{minipage}[t]{0.47\columnwidth}\raggedright
NOT modelled. Measured sensitivity: halving the board modulus moves the
driver ratio about 3 \%\strut
\end{minipage}\tabularnewline
\begin{minipage}[t]{0.47\columnwidth}\raggedright
assembly and reflow prestress\strut
\end{minipage} & \begin{minipage}[t]{0.47\columnwidth}\raggedright
NOT modelled. Measured: it moves the stabilised ratio by 0.11 \%\strut
\end{minipage}\tabularnewline
\begin{minipage}[t]{0.47\columnwidth}\raggedright
process residual stress\strut
\end{minipage} & \begin{minipage}[t]{0.47\columnwidth}\raggedright
NOT modelled, no data\strut
\end{minipage}\tabularnewline
\begin{minipage}[t]{0.47\columnwidth}\raggedright
interface delamination\strut
\end{minipage} & \begin{minipage}[t]{0.47\columnwidth}\raggedright
NOT modelled. Traction-separation parameters are unsourced\strut
\end{minipage}\tabularnewline
\begin{minipage}[t]{0.47\columnwidth}\raggedright
solder crack propagation\strut
\end{minipage} & \begin{minipage}[t]{0.47\columnwidth}\raggedright
NOT modelled. Not needed for ranking\strut
\end{minipage}\tabularnewline
\begin{minipage}[t]{0.47\columnwidth}\raggedright
\textbf{RDL and TSV cracking}\strut
\end{minipage} & \begin{minipage}[t]{0.47\columnwidth}\raggedright
\textbf{NOT modelled. The largest unmodelled mechanism that could change
ranking}\strut
\end{minipage}\tabularnewline
\begin{minipage}[t]{0.47\columnwidth}\raggedright
moisture, hygroscopic swelling\strut
\end{minipage} & \begin{minipage}[t]{0.47\columnwidth}\raggedright
OUT OF SCOPE, declared. The claim is dry steady operating ETM\strut
\end{minipage}\tabularnewline
\bottomrule
\end{longtable}

\newpage

\hypertarget{coupling-directions}{%
\section{15. Coupling directions}\label{coupling-directions}}

\begin{longtable}[]{@{}lll@{}}
\toprule
\begin{minipage}[b]{0.30\columnwidth}\raggedright
from -\textgreater{} to\strut
\end{minipage} & \begin{minipage}[b]{0.30\columnwidth}\raggedright
status\strut
\end{minipage} & \begin{minipage}[b]{0.30\columnwidth}\raggedright
measured magnitude\strut
\end{minipage}\tabularnewline
\midrule
\endhead
\begin{minipage}[t]{0.30\columnwidth}\raggedright
electrical -\textgreater{} thermal\strut
\end{minipage} & \begin{minipage}[t]{0.30\columnwidth}\raggedright
\textbf{solved, iterated}\strut
\end{minipage} & \begin{minipage}[t]{0.30\columnwidth}\raggedright
the whole loss\strut
\end{minipage}\tabularnewline
\begin{minipage}[t]{0.30\columnwidth}\raggedright
thermal -\textgreater{} electrical\strut
\end{minipage} & \begin{minipage}[t]{0.30\columnwidth}\raggedright
\textbf{solved, iterated}\strut
\end{minipage} & \begin{minipage}[t]{0.30\columnwidth}\raggedright
rho(T) moves R by 22.8 \% over the cycle\strut
\end{minipage}\tabularnewline
\begin{minipage}[t]{0.30\columnwidth}\raggedright
thermal -\textgreater{} mechanical\strut
\end{minipage} & \begin{minipage}[t]{0.30\columnwidth}\raggedright
\textbf{solved, one way}\strut
\end{minipage} & \begin{minipage}[t]{0.30\columnwidth}\raggedright
the whole thermal strain\strut
\end{minipage}\tabularnewline
\begin{minipage}[t]{0.30\columnwidth}\raggedright
geometry -\textgreater{} electrical / thermal / mechanical\strut
\end{minipage} & \begin{minipage}[t]{0.30\columnwidth}\raggedright
\textbf{solved}\strut
\end{minipage} & \begin{minipage}[t]{0.30\columnwidth}\raggedright
Sections 8, 12, 13\strut
\end{minipage}\tabularnewline
\begin{minipage}[t]{0.30\columnwidth}\raggedright
cooling -\textgreater{} electrical / thermal / mechanical\strut
\end{minipage} & \begin{minipage}[t]{0.30\columnwidth}\raggedright
\textbf{solved}\strut
\end{minipage} & \begin{minipage}[t]{0.30\columnwidth}\raggedright
Section 11, and Section 17\strut
\end{minipage}\tabularnewline
\begin{minipage}[t]{0.30\columnwidth}\raggedright
mechanical -\textgreater{} electrical\strut
\end{minipage} & \begin{minipage}[t]{0.30\columnwidth}\raggedright
\textbf{absent}\strut
\end{minipage} & \begin{minipage}[t]{0.30\columnwidth}\raggedright
0.16 \% on R, against 22.8 \% from rho(T) already in the loop\strut
\end{minipage}\tabularnewline
\begin{minipage}[t]{0.30\columnwidth}\raggedright
mechanical -\textgreater{} thermal\strut
\end{minipage} & \begin{minipage}[t]{0.30\columnwidth}\raggedright
\textbf{absent}\strut
\end{minipage} & \begin{minipage}[t]{0.30\columnwidth}\raggedright
1.5 K for a 4 \% TIM bond-line change\strut
\end{minipage}\tabularnewline
\begin{minipage}[t]{0.30\columnwidth}\raggedright
damage -\textgreater{} electrical\strut
\end{minipage} & \begin{minipage}[t]{0.30\columnwidth}\raggedright
\textbf{absent}\strut
\end{minipage} & \begin{minipage}[t]{0.30\columnwidth}\raggedright
would be +33 \% on R at 25 \% of a joint lost\strut
\end{minipage}\tabularnewline
\begin{minipage}[t]{0.30\columnwidth}\raggedright
damage -\textgreater{} thermal\strut
\end{minipage} & \begin{minipage}[t]{0.30\columnwidth}\raggedright
\textbf{absent}\strut
\end{minipage} & \begin{minipage}[t]{0.30\columnwidth}\raggedright
would be +12.5 K at 25 \% of the die delaminated\strut
\end{minipage}\tabularnewline
\begin{minipage}[t]{0.30\columnwidth}\raggedright
electromigration -\textgreater{} electrical\strut
\end{minipage} & \begin{minipage}[t]{0.30\columnwidth}\raggedright
\textbf{OPEN}\strut
\end{minipage} & \begin{minipage}[t]{0.30\columnwidth}\raggedright
the current-density allowable itself is unsourced\strut
\end{minipage}\tabularnewline
\bottomrule
\end{longtable}

\textbf{The framework is at ETM-L1}: electrical and thermal iterated to
a fixed point, thermal driving mechanics one way, mechanics returning
nothing.

That is a deliberate, measured choice and not an omission. The two
absent arrows before damage are 0.7 \% and 2.9 \% of terms already in
the loop, and both are common-mode across architectures, so solving them
iteratively would cost solver time to chase a term two orders of
magnitude smaller than one already present. \textbf{It is therefore
accurate to describe this framework as geometry-coupled
electrothermal-mechanical co-design with one-way mechanical response in
the intact regime, and inaccurate to call it fully coupled.}

ETM-L3, with damage feeding back into conductance and stiffness, would
be required to predict failure PROGRESSION. Nothing in this framework
currently claims to.

\newpage

\hypertarget{validation-matrix}{%
\section{16. Validation matrix}\label{validation-matrix}}

\begin{longtable}[]{@{}llll@{}}
\toprule
\begin{minipage}[b]{0.22\columnwidth}\raggedright
quantity\strut
\end{minipage} & \begin{minipage}[b]{0.22\columnwidth}\raggedright
method\strut
\end{minipage} & \begin{minipage}[b]{0.22\columnwidth}\raggedright
independent?\strut
\end{minipage} & \begin{minipage}[b]{0.22\columnwidth}\raggedright
result\strut
\end{minipage}\tabularnewline
\midrule
\endhead
\begin{minipage}[t]{0.22\columnwidth}\raggedright
RDL sheet PDE\strut
\end{minipage} & \begin{minipage}[t]{0.22\columnwidth}\raggedright
MOOSE\strut
\end{minipage} & \begin{minipage}[t]{0.22\columnwidth}\raggedright
yes\strut
\end{minipage} & \begin{minipage}[t]{0.22\columnwidth}\raggedright
0.33 to 0.45 \%\strut
\end{minipage}\tabularnewline
\begin{minipage}[t]{0.22\columnwidth}\raggedright
3-D conduction, uniform bar\strut
\end{minipage} & \begin{minipage}[t]{0.22\columnwidth}\raggedright
analytic\strut
\end{minipage} & \begin{minipage}[t]{0.22\columnwidth}\raggedright
yes\strut
\end{minipage} & \begin{minipage}[t]{0.22\columnwidth}\raggedright
exact at every mesh\strut
\end{minipage}\tabularnewline
\begin{minipage}[t]{0.22\columnwidth}\raggedright
3-D conduction, spreading\strut
\end{minipage} & \begin{minipage}[t]{0.22\columnwidth}\raggedright
analytic \texttt{rho/(4a)}\strut
\end{minipage} & \begin{minipage}[t]{0.22\columnwidth}\raggedright
yes\strut
\end{minipage} & \begin{minipage}[t]{0.22\columnwidth}\raggedright
+0.9 \% at 6 cells per radius\strut
\end{minipage}\tabularnewline
\begin{minipage}[t]{0.22\columnwidth}\raggedright
\textbf{transition resistance, 5 geometries}\strut
\end{minipage} & \begin{minipage}[t]{0.22\columnwidth}\raggedright
\textbf{MOOSE}\strut
\end{minipage} & \begin{minipage}[t]{0.22\columnwidth}\raggedright
\textbf{yes}\strut
\end{minipage} & \begin{minipage}[t]{0.22\columnwidth}\raggedright
\textbf{0.25 to 0.73 \%, 5/5 ACCEPTED}\strut
\end{minipage}\tabularnewline
\begin{minipage}[t]{0.22\columnwidth}\raggedright
TSV single-via DC resistance\strut
\end{minipage} & \begin{minipage}[t]{0.22\columnwidth}\raggedright
pre-existing framework, independent implementation\strut
\end{minipage} & \begin{minipage}[t]{0.22\columnwidth}\raggedright
yes\strut
\end{minipage} & \begin{minipage}[t]{0.22\columnwidth}\raggedright
exact\strut
\end{minipage}\tabularnewline
\begin{minipage}[t]{0.22\columnwidth}\raggedright
TSV skin ratio at 1 MHz\strut
\end{minipage} & \begin{minipage}[t]{0.22\columnwidth}\raggedright
Bessel solution\strut
\end{minipage} & \begin{minipage}[t]{0.22\columnwidth}\raggedright
yes\strut
\end{minipage} & \begin{minipage}[t]{0.22\columnwidth}\raggedright
delta/r = 4.41, ratio 1.0001\strut
\end{minipage}\tabularnewline
\begin{minipage}[t]{0.22\columnwidth}\raggedright
effective k\_z, all layers\strut
\end{minipage} & \begin{minipage}[t]{0.22\columnwidth}\raggedright
explicit-geometry RVE\strut
\end{minipage} & \begin{minipage}[t]{0.22\columnwidth}\raggedright
yes\strut
\end{minipage} & \begin{minipage}[t]{0.22\columnwidth}\raggedright
Voigt within 1 \%\strut
\end{minipage}\tabularnewline
\begin{minipage}[t]{0.22\columnwidth}\raggedright
effective k\_xy, bump layers\strut
\end{minipage} & \begin{minipage}[t]{0.22\columnwidth}\raggedright
explicit-geometry RVE\strut
\end{minipage} & \begin{minipage}[t]{0.22\columnwidth}\raggedright
yes\strut
\end{minipage} & \begin{minipage}[t]{0.22\columnwidth}\raggedright
Maxwell within 1.5 \%\strut
\end{minipage}\tabularnewline
\begin{minipage}[t]{0.22\columnwidth}\raggedright
effective k\_xy, TSV arrays\strut
\end{minipage} & \begin{minipage}[t]{0.22\columnwidth}\raggedright
explicit-geometry RVE\strut
\end{minipage} & \begin{minipage}[t]{0.22\columnwidth}\raggedright
yes\strut
\end{minipage} & \begin{minipage}[t]{0.22\columnwidth}\raggedright
Maxwell 11 to 39 \% high; the pre-existing model 19 to 74 \% low\strut
\end{minipage}\tabularnewline
\begin{minipage}[t]{0.22\columnwidth}\raggedright
unit-cell periodicity\strut
\end{minipage} & \begin{minipage}[t]{0.22\columnwidth}\raggedright
1x1 against 2x2\strut
\end{minipage} & \begin{minipage}[t]{0.22\columnwidth}\raggedright
yes\strut
\end{minipage} & \begin{minipage}[t]{0.22\columnwidth}\raggedright
0.000 \%\strut
\end{minipage}\tabularnewline
\begin{minipage}[t]{0.22\columnwidth}\raggedright
cooling fluid side\strut
\end{minipage} & \begin{minipage}[t]{0.22\columnwidth}\raggedright
analytic energy balance\strut
\end{minipage} & \begin{minipage}[t]{0.22\columnwidth}\raggedright
yes\strut
\end{minipage} & \begin{minipage}[t]{0.22\columnwidth}\raggedright
exact\strut
\end{minipage}\tabularnewline
\begin{minipage}[t]{0.22\columnwidth}\raggedright
\textbf{cooling solid side}\strut
\end{minipage} & \begin{minipage}[t]{0.22\columnwidth}\raggedright
\textbf{MOOSE, explicit fins}\strut
\end{minipage} & \begin{minipage}[t]{0.22\columnwidth}\raggedright
\textbf{yes}\strut
\end{minipage} & \begin{minipage}[t]{0.22\columnwidth}\raggedright
\textbf{3.78 \% at the production geometry; 10.11 \% for short thick
fins, REJECTED}\strut
\end{minipage}\tabularnewline
\begin{minipage}[t]{0.22\columnwidth}\raggedright
converter loss decomposition\strut
\end{minipage} & \begin{minipage}[t]{0.22\columnwidth}\raggedright
upstream aggregate\strut
\end{minipage} & \begin{minipage}[t]{0.22\columnwidth}\raggedright
yes\strut
\end{minipage} & \begin{minipage}[t]{0.22\columnwidth}\raggedright
machine precision, or the ledger refuses to total\strut
\end{minipage}\tabularnewline
\begin{minipage}[t]{0.22\columnwidth}\raggedright
spatial heat deposition\strut
\end{minipage} & \begin{minipage}[t]{0.22\columnwidth}\raggedright
electrical total\strut
\end{minipage} & \begin{minipage}[t]{0.22\columnwidth}\raggedright
yes\strut
\end{minipage} & \begin{minipage}[t]{0.22\columnwidth}\raggedright
within 1e-6 relative, guarded\strut
\end{minipage}\tabularnewline
\begin{minipage}[t]{0.22\columnwidth}\raggedright
solder PEEQ identity\strut
\end{minipage} & \begin{minipage}[t]{0.22\columnwidth}\raggedright
creep-only experiment\strut
\end{minipage} & \begin{minipage}[t]{0.22\columnwidth}\raggedright
yes\strut
\end{minipage} & \begin{minipage}[t]{0.22\columnwidth}\raggedright
1.0000\strut
\end{minipage}\tabularnewline
\begin{minipage}[t]{0.22\columnwidth}\raggedright
inelastic energy route\strut
\end{minipage} & \begin{minipage}[t]{0.22\columnwidth}\raggedright
solver mechanical strain\strut
\end{minipage} & \begin{minipage}[t]{0.22\columnwidth}\raggedright
yes\strut
\end{minipage} & \begin{minipage}[t]{0.22\columnwidth}\raggedright
0.79 to 1.56 \%, and 16 to 21x cleaner on the isochoric test\strut
\end{minipage}\tabularnewline
\begin{minipage}[t]{0.22\columnwidth}\raggedright
package thermal\strut
\end{minipage} & \begin{minipage}[t]{0.22\columnwidth}\raggedright
MOOSE, full 3-D\strut
\end{minipage} & \begin{minipage}[t]{0.22\columnwidth}\raggedright
yes\strut
\end{minipage} & \begin{minipage}[t]{0.22\columnwidth}\raggedright
ordering preserved\strut
\end{minipage}\tabularnewline
\begin{minipage}[t]{0.22\columnwidth}\raggedright
\textbf{package thermal against CFD}\strut
\end{minipage} & \begin{minipage}[t]{0.22\columnwidth}\raggedright
Icepak\strut
\end{minipage} & \begin{minipage}[t]{0.22\columnwidth}\raggedright
--\strut
\end{minipage} & \begin{minipage}[t]{0.22\columnwidth}\raggedright
\textbf{FAILED and discarded.} No CFD validation is claimed\strut
\end{minipage}\tabularnewline
\begin{minipage}[t]{0.22\columnwidth}\raggedright
\textbf{TSV loop AC resistance}\strut
\end{minipage} & \begin{minipage}[t]{0.22\columnwidth}\raggedright
HFSS\strut
\end{minipage} & \begin{minipage}[t]{0.22\columnwidth}\raggedright
--\strut
\end{minipage} & \begin{minipage}[t]{0.22\columnwidth}\raggedright
\textbf{ATTEMPTED, did not converge in budget}\strut
\end{minipage}\tabularnewline
\bottomrule
\end{longtable}

\newpage

\hypertarget{what-the-production-models-changed-measured}{%
\section{17. What the production models changed,
measured}\label{what-the-production-models-changed-measured}}

The same frozen comparison matrix was run twice, changing only which
models drive the path. BEFORE is the previous production path: one
global crowding constant, no intra-VR routing, a calibrated Robin
boundary with a uniform sink, a lumped VR. AFTER is the path described
in this document.

\begin{longtable}[]{@{}llllll@{}}
\toprule
architecture & mode & P\_elec & T\_max & die warpage & corner
shear\tabularnewline
\midrule
\endhead
Z1\_EMBEDDED & AFTER & 141.78 W & 123.53 C & 25.373 um & 4.2803
um\tabularnewline
Z1\_EMBEDDED & AFTER & 141.76 W & 123.53 C & 25.372 um & 4.2801
um\tabularnewline
Z1\_EMBEDDED & BEFORE & 125.10 W & 84.87 C & 20.042 um & 2.7447
um\tabularnewline
Z2\_BACKSIDE & AFTER & 143.61 W & 130.31 C & 22.622 um & 3.4101
um\tabularnewline
Z2\_BACKSIDE & AFTER & 143.60 W & 130.31 C & 22.621 um & 3.4100
um\tabularnewline
Z2\_BACKSIDE & AFTER & 136.25 W & 114.89 C & 18.825 um & 2.8842
um\tabularnewline
Z2\_BACKSIDE & AFTER & 146.89 W & 136.53 C & 24.242 um & 3.6258
um\tabularnewline
Z2\_BACKSIDE & BEFORE & 126.19 W & 83.31 C & 17.236 um & 2.0754
um\tabularnewline
Z4\_POWER\_DIE & AFTER & 134.74 W & 96.32 C & 35.946 um & 3.6681
um\tabularnewline
Z4\_POWER\_DIE & BEFORE & 123.64 W & 83.30 C & 30.138 um & 2.9744
um\tabularnewline
\bottomrule
\end{longtable}

\newpage

\hypertarget{numerical-convergence}{%
\section{18. Numerical convergence}\label{numerical-convergence}}

\begin{longtable}[]{@{}ll@{}}
\toprule
\begin{minipage}[b]{0.47\columnwidth}\raggedright
model\strut
\end{minipage} & \begin{minipage}[b]{0.47\columnwidth}\raggedright
convergence evidence\strut
\end{minipage}\tabularnewline
\midrule
\endhead
\begin{minipage}[t]{0.47\columnwidth}\raggedright
RDL sheet\strut
\end{minipage} & \begin{minipage}[t]{0.47\columnwidth}\raggedright
grid-converged; 121 to 201 nodes moves the answer 0.15 \%\strut
\end{minipage}\tabularnewline
\begin{minipage}[t]{0.47\columnwidth}\raggedright
3-D transition\strut
\end{minipage} & \begin{minipage}[t]{0.47\columnwidth}\raggedright
mesh drift 0.03 to 0.26 \% between refinement levels\strut
\end{minipage}\tabularnewline
\begin{minipage}[t]{0.47\columnwidth}\raggedright
RVE homogenization\strut
\end{minipage} & \begin{minipage}[t]{0.47\columnwidth}\raggedright
1x1 against 2x2 periodic cell, 0.000 \%\strut
\end{minipage}\tabularnewline
\begin{minipage}[t]{0.47\columnwidth}\raggedright
global thermal / structural\strut
\end{minipage} & \begin{minipage}[t]{0.47\columnwidth}\raggedright
element budget is a hard gate at the measured solver ceiling, never a
silent fallback\strut
\end{minipage}\tabularnewline
\begin{minipage}[t]{0.47\columnwidth}\raggedright
local solder\strut
\end{minipage} & \begin{minipage}[t]{0.47\columnwidth}\raggedright
stabilised cycle declared as a cycle-to-cycle change below 5 \%;
achieved 0.03 to 0.22 \%\strut
\end{minipage}\tabularnewline
\begin{minipage}[t]{0.47\columnwidth}\raggedright
energy integration\strut
\end{minipage} & \begin{minipage}[t]{0.47\columnwidth}\raggedright
tightened until the answer stopped moving; time step and output
frequency both refined\strut
\end{minipage}\tabularnewline
\begin{minipage}[t]{0.47\columnwidth}\raggedright
ET fixed point\strut
\end{minipage} & \begin{minipage}[t]{0.47\columnwidth}\raggedright
iterated on LOCAL layer temperature to below 0.25 K\strut
\end{minipage}\tabularnewline
\begin{minipage}[t]{0.47\columnwidth}\raggedright
cooling fluid coupling\strut
\end{minipage} & \begin{minipage}[t]{0.47\columnwidth}\raggedright
fluid temperature iterated to below 0.02 K\strut
\end{minipage}\tabularnewline
\begin{minipage}[t]{0.47\columnwidth}\raggedright
spatial heat\strut
\end{minipage} & \begin{minipage}[t]{0.47\columnwidth}\raggedright
deposited total checked against the electrical total, guarded at 1e-6
relative\strut
\end{minipage}\tabularnewline
\bottomrule
\end{longtable}

\newpage

\hypertarget{defects-found-and-corrected}{%
\section{19. Defects found and
corrected}\label{defects-found-and-corrected}}

This framework has a standing rule: a defect is recorded with the
mechanism that let it survive, not only with its fix. Fifty-three are
logged. The ones an auditor should know about:

\begin{longtable}[]{@{}llll@{}}
\toprule
\begin{minipage}[b]{0.22\columnwidth}\raggedright
id\strut
\end{minipage} & \begin{minipage}[b]{0.22\columnwidth}\raggedright
defect\strut
\end{minipage} & \begin{minipage}[b]{0.22\columnwidth}\raggedright
how it survived\strut
\end{minipage} & \begin{minipage}[b]{0.22\columnwidth}\raggedright
status\strut
\end{minipage}\tabularnewline
\midrule
\endhead
\begin{minipage}[t]{0.22\columnwidth}\raggedright
30\strut
\end{minipage} & \begin{minipage}[t]{0.22\columnwidth}\raggedright
vertical-tier losses summed per-regulator with system-wide RDL
loss\strut
\end{minipage} & \begin{minipage}[t]{0.22\columnwidth}\raggedright
no scope type existed\strut
\end{minipage} & \begin{minipage}[t]{0.22\columnwidth}\raggedright
FIXED, and made structural: mixing scopes now raises\strut
\end{minipage}\tabularnewline
\begin{minipage}[t]{0.22\columnwidth}\raggedright
40\strut
\end{minipage} & \begin{minipage}[t]{0.22\columnwidth}\raggedright
the creep law was fitted in a stress variable the solver does not
consume\strut
\end{minipage} & \begin{minipage}[t]{0.22\columnwidth}\raggedright
the fit and the solver were never compared in the same variable\strut
\end{minipage} & \begin{minipage}[t]{0.22\columnwidth}\raggedright
FIXED by removing the yield threshold\strut
\end{minipage}\tabularnewline
\begin{minipage}[t]{0.22\columnwidth}\raggedright
44\strut
\end{minipage} & \begin{minipage}[t]{0.22\columnwidth}\raggedright
stale current-facing claims propagated into both reports\strut
\end{minipage} & \begin{minipage}[t]{0.22\columnwidth}\raggedright
each sentence was true when written\strut
\end{minipage} & \begin{minipage}[t]{0.22\columnwidth}\raggedright
FIXED at source\strut
\end{minipage}\tabularnewline
\begin{minipage}[t]{0.22\columnwidth}\raggedright
46\strut
\end{minipage} & \begin{minipage}[t]{0.22\columnwidth}\raggedright
two parsers stood on opposite sides of a step boundary, inflating every
inelastic-energy number by 18.5 \%\strut
\end{minipage} & \begin{minipage}[t]{0.22\columnwidth}\raggedright
both parsers were individually correct\strut
\end{minipage} & \begin{minipage}[t]{0.22\columnwidth}\raggedright
FIXED, then ELIMINATED by reading the solver's own mechanical
strain\strut
\end{minipage}\tabularnewline
\begin{minipage}[t]{0.22\columnwidth}\raggedright
47\strut
\end{minipage} & \begin{minipage}[t]{0.22\columnwidth}\raggedright
the per-via current limit was a primitive in amperes, and it is the only
thing that sets field resistance\strut
\end{minipage} & \begin{minipage}[t]{0.22\columnwidth}\raggedright
it looked like a device property\strut
\end{minipage} & \begin{minipage}[t]{0.22\columnwidth}\raggedright
FIXED: \texttt{J\_design\ *\ A\_c}\strut
\end{minipage}\tabularnewline
\begin{minipage}[t]{0.22\columnwidth}\raggedright
48\strut
\end{minipage} & \begin{minipage}[t]{0.22\columnwidth}\raggedright
one global crowding constant for four interconnect classes\strut
\end{minipage} & \begin{minipage}[t]{0.22\columnwidth}\raggedright
it was measured once, on one geometry\strut
\end{minipage} & \begin{minipage}[t]{0.22\columnwidth}\raggedright
FIXED: solved per class, cross-validated\strut
\end{minipage}\tabularnewline
\begin{minipage}[t]{0.22\columnwidth}\raggedright
49\strut
\end{minipage} & \begin{minipage}[t]{0.22\columnwidth}\raggedright
the TSV liner was absent from the conductor area\strut
\end{minipage} & \begin{minipage}[t]{0.22\columnwidth}\raggedright
never modelled\strut
\end{minipage} & \begin{minipage}[t]{0.22\columnwidth}\raggedright
FIXED\strut
\end{minipage}\tabularnewline
\begin{minipage}[t]{0.22\columnwidth}\raggedright
50\strut
\end{minipage} & \begin{minipage}[t]{0.22\columnwidth}\raggedright
the converter area model assumes four-layer stacking and its own
docstring says the opposite\strut
\end{minipage} & \begin{minipage}[t]{0.22\columnwidth}\raggedright
\texttt{max()} reads like \texttt{sum()} at a glance\strut
\end{minipage} & \begin{minipage}[t]{0.22\columnwidth}\raggedright
DOCUMENTED; upstream is a separate repository\strut
\end{minipage}\tabularnewline
\begin{minipage}[t]{0.22\columnwidth}\raggedright
51\strut
\end{minipage} & \begin{minipage}[t]{0.22\columnwidth}\raggedright
intra-converter routing absent from every loss term\strut
\end{minipage} & \begin{minipage}[t]{0.22\columnwidth}\raggedright
the converter was a black box with a loss number\strut
\end{minipage} & \begin{minipage}[t]{0.22\columnwidth}\raggedright
FIXED: now in the total\strut
\end{minipage}\tabularnewline
\begin{minipage}[t]{0.22\columnwidth}\raggedright
52\strut
\end{minipage} & \begin{minipage}[t]{0.22\columnwidth}\raggedright
the calibrated cooling coefficient is unreachable by any laminar cooler
and hides 7.4 \% of the load in uncharged pumping\strut
\end{minipage} & \begin{minipage}[t]{0.22\columnwidth}\raggedright
it was calibrated to a temperature, and it hit it\strut
\end{minipage} & \begin{minipage}[t]{0.22\columnwidth}\raggedright
FIXED: derived from hardware\strut
\end{minipage}\tabularnewline
\begin{minipage}[t]{0.22\columnwidth}\raggedright
\textbf{53}\strut
\end{minipage} & \begin{minipage}[t]{0.22\columnwidth}\raggedright
\textbf{the current-status generator was internally consistent across
both reports and consistently stale}\strut
\end{minipage} & \begin{minipage}[t]{0.22\columnwidth}\raggedright
\textbf{cross-report comparison passes when both are stale in the same
way}\strut
\end{minipage} & \begin{minipage}[t]{0.22\columnwidth}\raggedright
\textbf{FIXED: a single authoritative contract, with mutation and
injection tests}\strut
\end{minipage}\tabularnewline
\bottomrule
\end{longtable}

Four faults in this phase's own new code were caught before they reached
a conclusion and are recorded for the same reason: a zero-width mesh
cell that produced a silently wrong resistance whose only symptom was a
current-balance residual; dead-ended traces that made a three-trace case
look 2.5x worse; a planar intra-VR layout that returned 61 \% of the
load because it assumed a topology the architecture does not have; and a
fit-quality metric that was misleading in both directions at once.

\newpage

\hypertarget{open-physical-gaps}{%
\section{20. Open physical gaps}\label{open-physical-gaps}}

\begin{longtable}[]{@{}lll@{}}
\toprule
\begin{minipage}[b]{0.30\columnwidth}\raggedright
gap\strut
\end{minipage} & \begin{minipage}[b]{0.30\columnwidth}\raggedright
why it is open\strut
\end{minipage} & \begin{minipage}[b]{0.30\columnwidth}\raggedright
consequence\strut
\end{minipage}\tabularnewline
\midrule
\endhead
\begin{minipage}[t]{0.30\columnwidth}\raggedright
\textbf{high-stress SAC305 constitutive response}\strut
\end{minipage} & \begin{minipage}[t]{0.30\columnwidth}\raggedright
no primary data above \textasciitilde{}16 MPa retrievable; four P0 and
five P1 papers behind institutional access\strut
\end{minipage} & \begin{minipage}[t]{0.30\columnwidth}\raggedright
\textbf{the dominant blocker.} Every absolute reliability number is
conditional on it\strut
\end{minipage}\tabularnewline
\begin{minipage}[t]{0.30\columnwidth}\raggedright
\textbf{required life}\strut
\end{minipage} & \begin{minipage}[t]{0.30\columnwidth}\raggedright
not published for the reference system, and inventing one is
forbidden\strut
\end{minipage} & \begin{minipage}[t]{0.30\columnwidth}\raggedright
the falsification gate cannot be evaluated\strut
\end{minipage}\tabularnewline
\begin{minipage}[t]{0.30\columnwidth}\raggedright
\textbf{mechanical allowable}\strut
\end{minipage} & \begin{minipage}[t]{0.30\columnwidth}\raggedright
nine candidates tested, none admissible\strut
\end{minipage} & \begin{minipage}[t]{0.30\columnwidth}\raggedright
mechanics ranks but cannot constrain\strut
\end{minipage}\tabularnewline
\begin{minipage}[t]{0.30\columnwidth}\raggedright
electromigration current-density allowable\strut
\end{minipage} & \begin{minipage}[t]{0.30\columnwidth}\raggedright
no primary study with its temperature, lifetime and geometry\strut
\end{minipage} & \begin{minipage}[t]{0.30\columnwidth}\raggedright
feasibility is a map over J, never a verdict\strut
\end{minipage}\tabularnewline
\begin{minipage}[t]{0.30\columnwidth}\raggedright
package clearances\strut
\end{minipage} & \begin{minipage}[t]{0.30\columnwidth}\raggedright
three unsourced\strut
\end{minipage} & \begin{minipage}[t]{0.30\columnwidth}\raggedright
no partition template may be called buildable\strut
\end{minipage}\tabularnewline
\begin{minipage}[t]{0.30\columnwidth}\raggedright
cohesive-zone parameters\strut
\end{minipage} & \begin{minipage}[t]{0.30\columnwidth}\raggedright
traction-separation laws unsourced\strut
\end{minipage} & \begin{minipage}[t]{0.30\columnwidth}\raggedright
delamination cannot be predicted\strut
\end{minipage}\tabularnewline
\begin{minipage}[t]{0.30\columnwidth}\raggedright
polymer viscoelastic spectra\strut
\end{minipage} & \begin{minipage}[t]{0.30\columnwidth}\raggedright
no DMA data for the specific laminate\strut
\end{minipage} & \begin{minipage}[t]{0.30\columnwidth}\raggedright
a modulus sensitivity is carried instead of a Prony series\strut
\end{minipage}\tabularnewline
\begin{minipage}[t]{0.30\columnwidth}\raggedright
Z4 power-die layer geometry\strut
\end{minipage} & \begin{minipage}[t]{0.30\columnwidth}\raggedright
the global mesh has no such layer\strut
\end{minipage} & \begin{minipage}[t]{0.30\columnwidth}\raggedright
the Z4 thermal result is conditional on depositing its heat in the
interface layer\strut
\end{minipage}\tabularnewline
\begin{minipage}[t]{0.30\columnwidth}\raggedright
conjugate CFD validation\strut
\end{minipage} & \begin{minipage}[t]{0.30\columnwidth}\raggedright
the earlier attempt returned a solver clamp\strut
\end{minipage} & \begin{minipage}[t]{0.30\columnwidth}\raggedright
no CFD validation is claimed\strut
\end{minipage}\tabularnewline
\bottomrule
\end{longtable}

\begin{figure}
\centering
\includegraphics{/home/beams0/MGHARIB/Mohamed_stuff/HVR_OPT/results/phase0mr/FIG_0MR_K_closure_ladder.png}
\caption{Closure ladder}
\end{figure}

\newpage

\hypertarget{which-claims-are-safe-and-which-are-not}{%
\section{21. Which claims are safe, and which are
not}\label{which-claims-are-safe-and-which-are-not}}

\hypertarget{safe-to-publish-now}{%
\subsection{21.1 Safe to publish now}\label{safe-to-publish-now}}

\begin{itemize}
\tightlist
\item
  Vertical placement of the converter changes both the electrical path
  and the local solder response, and the two trade against each other.
\item
  Moving conversion below the substrate puts the 1 V rail on the TSVs,
  and the electrical cost of that is quantified.
\item
  \textbf{Z2 carries a lower local fatigue-relevant inelastic response
  than Z1}, under every admitted constitutive model, every admitted life
  law, every profile and all three response metrics. The ordering
  originates upstream of the solder model. It is a RESPONSE ordering:
  calling it less damage would presume a sourced damage law, and that is
  OPEN.
\item
  The conditional relative-life envelope is \textbf{1.19 to 1.22} at the
  operating baseline.
\item
  Under minimum-capability provisioning the TSV conductor cross-section
  cancels out of the field resistance, and radius is a feasibility
  variable rather than a loss lever.
\item
  A single global current-crowding multiplier is the wrong functional
  form; the measured spread over 107 geometries is 2.4x.
\item
  Intra-converter metallization is a real loss term of order 0.4 to 3.9
  \% of the load and was previously absent.
\item
  The calibrated cooling coefficient is not reachable by any laminar
  microchannel architecture, and reaching it costs about 7.4 \% of the
  load in pumping power.
\item
  A single-level VR does not fit the die shadow at 1 kW; vertical
  partitioning is a precondition.
\item
  What CalculiX integrates in a \texttt{*VISCO} step is an overstress
  law with no rate-independent branch, traced in the solver source.
\end{itemize}

\hypertarget{not-safe-and-why}{%
\subsection{21.2 NOT safe, and why}\label{not-safe-and-why}}

\begin{longtable}[]{@{}ll@{}}
\toprule
\begin{minipage}[b]{0.47\columnwidth}\raggedright
claim\strut
\end{minipage} & \begin{minipage}[b]{0.47\columnwidth}\raggedright
why not\strut
\end{minipage}\tabularnewline
\midrule
\endhead
\begin{minipage}[t]{0.47\columnwidth}\raggedright
``Z2 lasts N cycles''\strut
\end{minipage} & \begin{minipage}[t]{0.47\columnwidth}\raggedright
no absolute life is defensible\strut
\end{minipage}\tabularnewline
\begin{minipage}[t]{0.47\columnwidth}\raggedright
``validated relative life''\strut
\end{minipage} & \begin{minipage}[t]{0.47\columnwidth}\raggedright
the magnitude is a conditional envelope; the wording is not earned\strut
\end{minipage}\tabularnewline
\begin{minipage}[t]{0.47\columnwidth}\raggedright
``the solder model is valid for this package''\strut
\end{minipage} & \begin{minipage}[t]{0.47\columnwidth}\raggedright
it is evaluated far outside its fitted stress window\strut
\end{minipage}\tabularnewline
\begin{minipage}[t]{0.47\columnwidth}\raggedright
``the creep law is SAC305''\strut
\end{minipage} & \begin{minipage}[t]{0.47\columnwidth}\raggedright
its data is Sn3.9Ag0.6Cu\strut
\end{minipage}\tabularnewline
\begin{minipage}[t]{0.47\columnwidth}\raggedright
``the stress extrapolation is common mode and cancels''\strut
\end{minipage} & \begin{minipage}[t]{0.47\columnwidth}\raggedright
measured, and it only partly cancels\strut
\end{minipage}\tabularnewline
\begin{minipage}[t]{0.47\columnwidth}\raggedright
``mechanics constrains the design''\strut
\end{minipage} & \begin{minipage}[t]{0.47\columnwidth}\raggedright
no admissible allowable exists\strut
\end{minipage}\tabularnewline
\begin{minipage}[t]{0.47\columnwidth}\raggedright
``the framework is fully ETM coupled''\strut
\end{minipage} & \begin{minipage}[t]{0.47\columnwidth}\raggedright
mechanics returns nothing; it is ETM-L1\strut
\end{minipage}\tabularnewline
\begin{minipage}[t]{0.47\columnwidth}\raggedright
``2.5D is eliminated''\strut
\end{minipage} & \begin{minipage}[t]{0.47\columnwidth}\raggedright
only the one-sided low-voltage lateral topology is\strut
\end{minipage}\tabularnewline
\begin{minipage}[t]{0.47\columnwidth}\raggedright
``TSV radius is irrelevant''\strut
\end{minipage} & \begin{minipage}[t]{0.47\columnwidth}\raggedright
it is not a field-loss lever; it is decisive for feasibility\strut
\end{minipage}\tabularnewline
\begin{minipage}[t]{0.47\columnwidth}\raggedright
``the thermal model is validated against CFD''\strut
\end{minipage} & \begin{minipage}[t]{0.47\columnwidth}\raggedright
the CFD attempt failed and was discarded\strut
\end{minipage}\tabularnewline
\begin{minipage}[t]{0.47\columnwidth}\raggedright
``this partition template is buildable''\strut
\end{minipage} & \begin{minipage}[t]{0.47\columnwidth}\raggedright
three package clearances are unsourced\strut
\end{minipage}\tabularnewline
\begin{minipage}[t]{0.47\columnwidth}\raggedright
``there is plastic and creep strain in the joint''\strut
\end{minipage} & \begin{minipage}[t]{0.47\columnwidth}\raggedright
there is one viscoplastic strain\strut
\end{minipage}\tabularnewline
\bottomrule
\end{longtable}

\newpage

\hypertarget{how-this-document-stays-true}{%
\section{22. How this document stays
true}\label{how-this-document-stays-true}}

Both this report and its chronological companion are generated from one
file. That meant, for eight phases, that comparing them against each
other could pass while both were stale --- and it did. The section
headed ``current scientific status'' printed a current verdict above six
bullets from a much earlier phase, and cross-report consistency was 25
of 25 the whole time.

The fix is the contract in Section 1.1. Every current-facing statement
derives from it, and three tests keep it honest:

\begin{itemize}
\tightlist
\item
  \textbf{mutation} --- change one contract field and both reports fail
  until regenerated;
\item
  \textbf{injection} --- each of five REAL stale sentences, copied
  verbatim from the artifact that contained it, is inserted into a
  current-facing block and must be caught;
\item
  \textbf{exoneration} --- a defect report must still be able to quote
  the defect it fixed, or the guard becomes unusable and someone will
  switch it off.
\end{itemize}

Every evidence path in the contract is checked to exist on disk. That
check rejected the first draft.

\newpage

\hypertarget{what-should-happen-next}{%
\section{23. What should happen next}\label{what-should-happen-next}}

In priority order, with the reason each is ahead of the next:

\begin{enumerate}
\def\labelenumi{\arabic{enumi}.}
\tightlist
\item
  \textbf{Characterise SAC305 above 16 MPa}, by coupon test and by
  acquiring the four P0 papers. It is the dominant blocker on every
  absolute number, and no amount of modelling removes it.
\item
  \textbf{Obtain a required life} from the system owner. Without it the
  falsification gate cannot be evaluated, whatever the mechanics says.
\item
  \textbf{Put the RDL into the production mechanical and thermal stack
  by default.} The ladder is built and measured; it moves the absolute
  response by up to 7 \% and the peak temperature by up to 4.7 K.
\item
  \textbf{Run a canonical single-channel Icepak benchmark} against the
  analytical correlation. The licence is available. Until it reproduces
  the correlation, no CFD validation may be claimed.
\item
  \textbf{Source the package clearances}, without which no partition
  template can be called buildable.
\item
  \textbf{Then, and only then, consider an optimizer.} Its preconditions
  are listed in Section 24 and none of them is currently met in full.
\end{enumerate}

\newpage

\hypertarget{the-optimizer-gate}{%
\section{24. The optimizer gate}\label{the-optimizer-gate}}

A global optimizer is \textbf{FORBIDDEN} and remains so. Its
preconditions, and where they stand:

\begin{longtable}[]{@{}lll@{}}
\toprule
\begin{minipage}[b]{0.30\columnwidth}\raggedright
precondition\strut
\end{minipage} & \begin{minipage}[b]{0.30\columnwidth}\raggedright
status after Phase 0M-R3C\strut
\end{minipage} & \begin{minipage}[b]{0.30\columnwidth}\raggedright
evidence\strut
\end{minipage}\tabularnewline
\midrule
\endhead
\begin{minipage}[t]{0.30\columnwidth}\raggedright
interconnect geometry and crowding validated\strut
\end{minipage} & \begin{minipage}[t]{0.30\columnwidth}\raggedright
\textbf{PARTIAL}\strut
\end{minipage} & \begin{minipage}[t]{0.30\columnwidth}\raggedright
Delta\_R cross-validated against MOOSE at 0.25-0.73 \%; TSV span
architecture-aware and consistent across all three domains. The EM
allowable is still OPEN.\strut
\end{minipage}\tabularnewline
\begin{minipage}[t]{0.30\columnwidth}\raggedright
routing scope physically correct\strut
\end{minipage} & \begin{minipage}[t]{0.30\columnwidth}\raggedright
\textbf{PARTIAL}\strut
\end{minipage} & \begin{minipage}[t]{0.30\columnwidth}\raggedright
Vertical spans are measured tier mid-plane to tier mid-plane and
overstate the build-up crossed by 1.63x to 65.55x (DEFECT \#62). The
correction is available behind PHASE0MR3B\_PHYSICAL\_SPANS=1 and is
default off, so it is known and not yet adopted.\strut
\end{minipage}\tabularnewline
\begin{minipage}[t]{0.30\columnwidth}\raggedright
cooling boundary tied to physical hardware\strut
\end{minipage} & \begin{minipage}[t]{0.30\columnwidth}\raggedright
\textbf{PARTIAL}\strut
\end{minipage} & \begin{minipage}[t]{0.30\columnwidth}\raggedright
the production point is inside a SOURCE\_BACKED correlation domain and
the solid side is MOOSE-validated at 3.78 \%. The fluid side has no
conjugate CFD validation, and 2300 \textless{} Re \textless{} 3000 has
no correlation at all.\strut
\end{minipage}\tabularnewline
\begin{minipage}[t]{0.30\columnwidth}\raggedright
architecture templates realizable\strut
\end{minipage} & \begin{minipage}[t]{0.30\columnwidth}\raggedright
\textbf{PARTIAL}\strut
\end{minipage} & \begin{minipage}[t]{0.30\columnwidth}\raggedright
UPGRADED from NO by 0M-R3C. A permutation of the existing components
(T1C\_REORDERED) closes the geometry for one Z1 and one Z2: worst
build-up aspect ratio 0.833, worst framework routing ratio 0.633. It
cannot reach YES while the microvia current allowable, the capture pad,
a vendor design rule and two package clearances are OPEN.\strut
\end{minipage}\tabularnewline
\begin{minipage}[t]{0.30\columnwidth}\raggedright
architecture geometry physically represented\strut
\end{minipage} & \begin{minipage}[t]{0.30\columnwidth}\raggedright
\textbf{PARTIAL}\strut
\end{minipage} & \begin{minipage}[t]{0.30\columnwidth}\raggedright
ARCH\_Z4's power die is explicit geometry but its construction is
unsourced, so its ranking stays conditional.\strut
\end{minipage}\tabularnewline
\begin{minipage}[t]{0.30\columnwidth}\raggedright
process margin on the critical interconnect\strut
\end{minipage} & \begin{minipage}[t]{0.30\columnwidth}\raggedright
\textbf{NO}\strut
\end{minipage} & \begin{minipage}[t]{0.30\columnwidth}\raggedright
ARCH\_Z2's TSV field sits at AR = 10.000, M\_AR = 1.000, and fails a
modest tolerance bound. 0M-R3A re-derived the sweep: the smallest
passing diameter is 32 um, not 34 um, and enlarging the via REDUCES the
required count. Neither is adopted into the frozen pair.\strut
\end{minipage}\tabularnewline
\begin{minipage}[t]{0.30\columnwidth}\raggedright
intra-VR via technology sourced\strut
\end{minipage} & \begin{minipage}[t]{0.30\columnwidth}\raggedright
\textbf{NO}\strut
\end{minipage} & \begin{minipage}[t]{0.30\columnwidth}\raggedright
The interconnect CLASS is now sourced (stacked build-up microvias
through embedded-component substrate power layers, Ostmann 2008 + Cauwe
2023), but five of twelve geometric parameters remain OPEN: microvia
diameter, capture pad, landing area, via current limit and per-segment
aspect-ratio limit.\strut
\end{minipage}\tabularnewline
\begin{minipage}[t]{0.30\columnwidth}\raggedright
dominant mechanical physics closed\strut
\end{minipage} & \begin{minipage}[t]{0.30\columnwidth}\raggedright
\textbf{NO}\strut
\end{minipage} & \begin{minipage}[t]{0.30\columnwidth}\raggedright
high-stress SAC305 constitutive response is CONDITIONAL; cohesive-zone
parameters OPEN; 6 of 8 enumerated failure mechanisms uncomputed.\strut
\end{minipage}\tabularnewline
\begin{minipage}[t]{0.30\columnwidth}\raggedright
an admissible mechanical criterion exists\strut
\end{minipage} & \begin{minipage}[t]{0.30\columnwidth}\raggedright
\textbf{NO}\strut
\end{minipage} & \begin{minipage}[t]{0.30\columnwidth}\raggedright
nine candidates tested against six admissibility conditions; none
passed.\strut
\end{minipage}\tabularnewline
\begin{minipage}[t]{0.30\columnwidth}\raggedright
a required life exists\strut
\end{minipage} & \begin{minipage}[t]{0.30\columnwidth}\raggedright
\textbf{NO}\strut
\end{minipage} & \begin{minipage}[t]{0.30\columnwidth}\raggedright
N\_required is OPEN.\strut
\end{minipage}\tabularnewline
\begin{minipage}[t]{0.30\columnwidth}\raggedright
absolute fatigue prediction validated\strut
\end{minipage} & \begin{minipage}[t]{0.30\columnwidth}\raggedright
\textbf{NO}\strut
\end{minipage} & \begin{minipage}[t]{0.30\columnwidth}\raggedright
the strain and energy formulations disagree 4.1x on this joint.\strut
\end{minipage}\tabularnewline
\bottomrule
\end{longtable}

Phase 0M-R3A DOWNGRADED the routing precondition from YES to PARTIAL,
because the repair that makes both candidates pass sits in an unsourced
via process window, and added a new precondition for whether that via
technology is sourced at all. Phase 0M-R2 had DOWNGRADED the cooling
precondition from YES to PARTIAL --- the model has no validated
correlation between Re 2300 and Re 3000, and an optimizer would be free
to walk into that hole --- and added a new one for whether the
architecture geometry is physically represented at all. \textbf{6 remain
NO, and any one of them is sufficient to keep the gate shut.}

\newpage

\hypertarget{reproducibility}{%
\section{25. Reproducibility}\label{reproducibility}}

\begin{longtable}[]{@{}ll@{}}
\toprule
\begin{minipage}[b]{0.47\columnwidth}\raggedright
item\strut
\end{minipage} & \begin{minipage}[b]{0.47\columnwidth}\raggedright
location\strut
\end{minipage}\tabularnewline
\midrule
\endhead
\begin{minipage}[t]{0.47\columnwidth}\raggedright
single source of truth for current state\strut
\end{minipage} & \begin{minipage}[t]{0.47\columnwidth}\raggedright
\texttt{results/CURRENT\_SCIENTIFIC\_STATUS.yaml::current\_state\_contract}\strut
\end{minipage}\tabularnewline
\begin{minipage}[t]{0.47\columnwidth}\raggedright
production interconnect model\strut
\end{minipage} & \begin{minipage}[t]{0.47\columnwidth}\raggedright
\texttt{scripts/phase0mr/interconnect\_production.py}\strut
\end{minipage}\tabularnewline
\begin{minipage}[t]{0.47\columnwidth}\raggedright
3-D conduction solver\strut
\end{minipage} & \begin{minipage}[t]{0.47\columnwidth}\raggedright
\texttt{scripts/phase0m/conduction3d.py}\strut
\end{minipage}\tabularnewline
\begin{minipage}[t]{0.47\columnwidth}\raggedright
independent cross-check\strut
\end{minipage} & \begin{minipage}[t]{0.47\columnwidth}\raggedright
\texttt{scripts/phase0mr/moose\_conduction.py}\strut
\end{minipage}\tabularnewline
\begin{minipage}[t]{0.47\columnwidth}\raggedright
intra-VR routing\strut
\end{minipage} & \begin{minipage}[t]{0.47\columnwidth}\raggedright
\texttt{scripts/phase0mr/routing\_production.py}\strut
\end{minipage}\tabularnewline
\begin{minipage}[t]{0.47\columnwidth}\raggedright
loss ledger, one owner per watt\strut
\end{minipage} & \begin{minipage}[t]{0.47\columnwidth}\raggedright
\texttt{scripts/phase0mr/ledger.py}\strut
\end{minipage}\tabularnewline
\begin{minipage}[t]{0.47\columnwidth}\raggedright
physical cooling\strut
\end{minipage} & \begin{minipage}[t]{0.47\columnwidth}\raggedright
\texttt{scripts/phase0mr/cooling\_production.py}\strut
\end{minipage}\tabularnewline
\begin{minipage}[t]{0.47\columnwidth}\raggedright
production ET fixed point\strut
\end{minipage} & \begin{minipage}[t]{0.47\columnwidth}\raggedright
\texttt{scripts/phase0mr/production\_et.py}\strut
\end{minipage}\tabularnewline
\begin{minipage}[t]{0.47\columnwidth}\raggedright
controlled replay\strut
\end{minipage} & \begin{minipage}[t]{0.47\columnwidth}\raggedright
\texttt{scripts/phase0mr/run\_replay.py}\strut
\end{minipage}\tabularnewline
\begin{minipage}[t]{0.47\columnwidth}\raggedright
current-state guard\strut
\end{minipage} & \begin{minipage}[t]{0.47\columnwidth}\raggedright
\texttt{scripts/phase0mr/defect53\_audit.py}\strut
\end{minipage}\tabularnewline
\begin{minipage}[t]{0.47\columnwidth}\raggedright
tests\strut
\end{minipage} & \begin{minipage}[t]{0.47\columnwidth}\raggedright
\texttt{tests/}, run with \texttt{python\ -m\ pytest\ tests\ -q}\strut
\end{minipage}\tabularnewline
\bottomrule
\end{longtable}

The repository is \textbf{not under version control}, so content
addressing is used: every authoritative artifact's md5 is recorded at
each phase boundary and the previous PDFs are archived byte-identically
under \texttt{results/archive/}.

Environment: Python 3.12.14, CalculiX 2.23, MOOSE heat\_transfer
\texttt{09515f92aa}, gmsh 4.15.2, AEDT 2026 R1 (HFSS and Icepak
licensed, Q3D not), pandoc 2.0.6 with xelatex.

\hypertarget{auditor-checklist}{%
\section{26. Auditor checklist}\label{auditor-checklist}}

Ten questions worth asking of this framework, with where to find the
answer:

\begin{longtable}[]{@{}ll@{}}
\toprule
\begin{minipage}[b]{0.47\columnwidth}\raggedright
question\strut
\end{minipage} & \begin{minipage}[b]{0.47\columnwidth}\raggedright
where\strut
\end{minipage}\tabularnewline
\midrule
\endhead
\begin{minipage}[t]{0.47\columnwidth}\raggedright
is any resistance a primitive?\strut
\end{minipage} & \begin{minipage}[t]{0.47\columnwidth}\raggedright
Section 8 --- no, and the classifier raises on anything
unclassified\strut
\end{minipage}\tabularnewline
\begin{minipage}[t]{0.47\columnwidth}\raggedright
can any loss term be double counted?\strut
\end{minipage} & \begin{minipage}[t]{0.47\columnwidth}\raggedright
Section 10 --- the ledger enforces one owner per watt and refuses to
total otherwise\strut
\end{minipage}\tabularnewline
\begin{minipage}[t]{0.47\columnwidth}\raggedright
is the cooling boundary fitted to its own answer?\strut
\end{minipage} & \begin{minipage}[t]{0.47\columnwidth}\raggedright
Section 11 --- no longer; it is derived from channel geometry and
flow\strut
\end{minipage}\tabularnewline
\begin{minipage}[t]{0.47\columnwidth}\raggedright
is any solver used outside what it can pose?\strut
\end{minipage} & \begin{minipage}[t]{0.47\columnwidth}\raggedright
Section 5 --- HFSS is explicitly NOT used for DC conduction, and
why\strut
\end{minipage}\tabularnewline
\begin{minipage}[t]{0.47\columnwidth}\raggedright
is any effective property unvalidated?\strut
\end{minipage} & \begin{minipage}[t]{0.47\columnwidth}\raggedright
Section 12 --- every tensor has an RVE error and a verdict\strut
\end{minipage}\tabularnewline
\begin{minipage}[t]{0.47\columnwidth}\raggedright
is the metric mesh-converged?\strut
\end{minipage} & \begin{minipage}[t]{0.47\columnwidth}\raggedright
Section 18\strut
\end{minipage}\tabularnewline
\begin{minipage}[t]{0.47\columnwidth}\raggedright
is any claim stale?\strut
\end{minipage} & \begin{minipage}[t]{0.47\columnwidth}\raggedright
Section 22 --- generated from a contract, with mutation and injection
tests\strut
\end{minipage}\tabularnewline
\begin{minipage}[t]{0.47\columnwidth}\raggedright
what would change the ranking?\strut
\end{minipage} & \begin{minipage}[t]{0.47\columnwidth}\raggedright
Section 20 --- RDL and TSV cracking is the largest unmodelled
candidate\strut
\end{minipage}\tabularnewline
\begin{minipage}[t]{0.47\columnwidth}\raggedright
what is conditional and on what?\strut
\end{minipage} & \begin{minipage}[t]{0.47\columnwidth}\raggedright
Section 14.2, and the uncertainty ledger\strut
\end{minipage}\tabularnewline
\begin{minipage}[t]{0.47\columnwidth}\raggedright
may an optimizer be built?\strut
\end{minipage} & \begin{minipage}[t]{0.47\columnwidth}\raggedright
Section 24 --- no, four preconditions remain unmet\strut
\end{minipage}\tabularnewline
\bottomrule
\end{longtable}

\hypertarget{symbols}{%
\section{27. Symbols}\label{symbols}}

\begin{longtable}[]{@{}lll@{}}
\toprule
symbol & meaning & units\tabularnewline
\midrule
\endhead
\texttt{rho(T)} & resistivity & Ohm.m\tabularnewline
\texttt{A\_c} & conductor cross-section & m2\tabularnewline
\texttt{J\_design} & provisioning current density & A/m2\tabularnewline
\texttt{J\_process\_max} & manufacturing capability bound &
A/m2\tabularnewline
\texttt{J\_reliability\_allow} & electromigration allowable ---
\textbf{OPEN} & A/m2\tabularnewline
\texttt{n\_req}, \texttt{n\_avail}, \texttt{n\_alloc} & required,
available and allocated via counts & --\tabularnewline
\texttt{dR\_transition} & 3-D resistance minus the 1-D reference &
Ohm\tabularnewline
\texttt{sigma\_s} & sheet conductance & S\tabularnewline
\texttt{k\_x,\ k\_y,\ k\_z} & effective conductivity tensor diagonal &
W/m.K\tabularnewline
\texttt{h(x)} & local film coefficient & W/m2.K\tabularnewline
\texttt{T\_f(x)} & local coolant temperature & C\tabularnewline
\texttt{eta\_fin} & fin efficiency & --\tabularnewline
\texttt{du\_corner} & corner-BGA relative in-plane displacement &
um\tabularnewline
\texttt{PEEQ} & accumulated equivalent inelastic strain &
--\tabularnewline
\texttt{ME} & solver mechanical strain, total minus thermal &
--\tabularnewline
\texttt{D} & per-cycle local inelastic RESPONSE metric (dPEEQ) &
--\tabularnewline
\texttt{N\_required} & required life --- \textbf{OPEN} &
cycles\tabularnewline
\bottomrule
\end{longtable}

\hypertarget{glossary}{%
\section{28. Glossary}\label{glossary}}

\textbf{DVPD} --- distributed vertical power delivery: converting close
to the load so the high-current rail travels vertically and briefly
rather than laterally.

\textbf{Overstress law} --- a creep law driven by the amount by which
the stress exceeds a yield surface, rather than by the stress itself. It
returns exactly zero below yield. This is what CalculiX integrates in a
\texttt{*VISCO} step.

\textbf{Homogenization} --- replacing a geometrically detailed layer
with an equivalent uniform material. Valid only with a stated error
against the explicit geometry.

\textbf{RVE} --- representative volume element: the smallest repeating
cell whose solution gives the effective property of the whole layer.

\textbf{Constriction / spreading resistance} --- the extra resistance
where current converges into or diverges out of a smaller cross-section.
Not captured by \texttt{rho\ L\ /\ A}.

\textbf{Fin efficiency} --- the fraction of a cooling fin's surface that
is as effective as its base. It falls as the fin gets taller or thinner.

\textbf{Caloric resistance} --- the temperature rise caused simply by
the coolant absorbing heat, \texttt{1/(2\ m\_dot\ c\_p)}. It is a
property of the flow, not of the surface, and a single film coefficient
hides it.

\textbf{Common mode} --- an error that multiplies both architectures by
the same factor and therefore cancels in their ratio. Whether it does
must be tested, not assumed.

\textbf{ETM-L1 / L2 / L3} --- coupling levels: L1 is electrical-thermal
iterated with one-way mechanics; L2 adds mechanical feedback into
conductance or geometry; L3 adds damage feeding back. This framework is
L1.

\end{document}
